\documentclass[11pt]{article}
\usepackage[T1]{fontenc}
\usepackage[utf8]{inputenc}
\usepackage{lmodern}
\usepackage[margin=1in]{geometry}
\usepackage{microtype}
\usepackage{amsmath,amssymb,amsthm,mathtools}
\usepackage{booktabs,longtable,array,tabularx}
\usepackage{enumitem}
\usepackage{xcolor}
\usepackage{listings}
\usepackage{hyperref}
\usepackage[nameinlink,noabbrev]{cleveref}
\usepackage{fancyhdr}
\usepackage{titlesec}
\usepackage{etoolbox}
\usepackage{url}

\hypersetup{
  colorlinks=true,
  linkcolor=blue!50!black,
  citecolor=blue!50!black,
  urlcolor=blue!55!black,
  pdftitle={FINER-iPRA: Finite-Interaction Normal Form with Explicit Residuals},
  pdfauthor={OpenAI, prepared for Vladimir Reshetnikov}
}

\definecolor{codebg}{RGB}{247,247,247}
\definecolor{statusgreen}{RGB}{28,110,70}
\definecolor{statusamber}{RGB}{160,96,0}
\definecolor{statusred}{RGB}{145,35,35}
\lstset{
  basicstyle=\ttfamily\small,
  backgroundcolor=\color{codebg},
  frame=single,
  breaklines=true,
  columns=fullflexible,
  keepspaces=true,
  showstringspaces=false,
  xleftmargin=0.5em,
  xrightmargin=0.5em
}

\newtheorem{theorem}{Theorem}[section]
\newtheorem{lemma}[theorem]{Lemma}
\newtheorem{proposition}[theorem]{Proposition}
\newtheorem{corollary}[theorem]{Corollary}
\newtheorem{definition}[theorem]{Definition}
\newtheorem{principle}[theorem]{Principle}
\newtheorem{barrier}[theorem]{Barrier}
\theoremstyle{remark}
\newtheorem{remark}[theorem]{Remark}
\newtheorem{example}[theorem]{Example}
\newtheorem{status}[theorem]{Status}

\newcommand{\iPRA}{\ensuremath{\mathsf{iPRA}}}
\newcommand{\PRAAx}{\ensuremath{\mathsf{PRAAx}}}
\newcommand{\Prf}{\ensuremath{\mathsf{Prf}}}
\newcommand{\Nat}{\ensuremath{\mathbb N}}
\newcommand{\Bool}{\ensuremath{\mathbb B}}
\newcommand{\True}{\ensuremath{\mathbf 1}}
\newcommand{\False}{\ensuremath{\mathbf 0}}
\newcommand{\seq}[3][]{\ensuremath{#2\vdash_{\PRAAx}^{#1}#3}}
\newcommand{\objseq}[2]{\ensuremath{#1\;\vdash[\PRAAx]\;#2}}
\newcommand{\FINER}{\textsc{Finer--iPRA}}
\newcommand{\CertCheck}{\ensuremath{\mathsf{Check}}}
\newcommand{\Obs}{\ensuremath{\mathsf{Obs}}}
\newcommand{\Sel}{\ensuremath{\mathsf{Sel}}}
\newcommand{\rank}{\ensuremath{\mathsf{rk}}}
\newcommand{\code}{\ensuremath{\mathsf{code}}}
\newcommand{\decode}{\ensuremath{\mathsf{decode}}}
\newcommand{\Resid}{\ensuremath{\mathsf{Resid}}}
\newcommand{\qf}{\ensuremath{\mathsf{QF}}}
\newcommand{\ok}{\textcolor{statusgreen}{\textbf{checked/inherited}}}
\newcommand{\human}{\textcolor{statusamber}{\textbf{human proof}}}
\newcommand{\outline}{\textcolor{statusamber}{\textbf{Lean outline}}}
\newcommand{\openstatus}{\textcolor{statusred}{\textbf{residual/open}}}

\pagestyle{fancy}
\fancyhf{}
\lhead{\FINER}
\rhead{Finite mathematics in \iPRA}
\cfoot{\thepage}
\setlength{\headheight}{14pt}

\title{\bfseries FINER--iPRA\\[0.35em]
\Large Finite-Interaction Normal Form with Explicit Residuals\\[0.5em]
\large A general framework for compiling mathematical proofs into intuitionistic primitive recursive arithmetic}
\author{Prepared for Vladimir Reshetnikov\\OpenAI}
\date{16 November 2026\\Repository baseline: \texttt{64bc0e997df6b8dbdf43f5ac9b3d9f74b4047ee5}}

\begin{document}
\maketitle

\begin{abstract}
The ProveIt repository contains a mature kernel for ordinary intuitionistic primitive recursive arithmetic, a substantial checked library of object-theory arithmetic derivations, and a much larger collection of research packets that recast analysis, algebra, number theory, geometry, probability, and topology in finitary form.  The remaining difficulty is no longer a lack of candidate encodings.  It is the absence of one implementation discipline that turns a theorem-specific finitary argument into a small, auditable value of the literal Lean type \(\Prf\;\PRAAx\;\Gamma\;\varphi\), while preserving the exact boundary between verification and the production of deep quantitative data.

This report proposes \FINER, ``Finite-Interaction Normal form with Explicit Residuals.''  A mathematical assertion is first specialized to a finite query.  Its evidence is packed into one bounded master certificate, checked by a fixed primitive-recursive Boolean program, and selected by the already checked bounded-search development.  Deterministic computation is reflected by one quantifier-free induction; non-unique witnesses are handled by backward viability rather than hidden choice; outputs descend only through exact, quotient, or metric observers; and every unavailable height bound, separation margin, modulus, finite cover, or analytic constant remains as an explicit formula in the final sequent.  Supplying object-language proofs of those residual formulas closes the theorem by the checked simultaneous-cut rule.

Two additions sharpen earlier ProveIt frameworks.  First, a finite monotone demand-closure theorem handles genuinely cyclic precision and resource dependencies: iteration of \(d\mapsto d\vee F(d)\) on a finite product of chains reaches the least prefixed point after at most the sum of the coordinate capacities.  Its proof uses a primitive-recursive rank and quantifier-free induction.  Second, finite observational descent replaces unnecessary global gluing: edge compatibility on a supplied spanning tree gives exact equality, quotient agreement, or an explicit accumulated metric error at a root observation, without constructing an infinitary quotient or selecting a global witness.

The report gives complete human proofs of the reusable schemas, explains what they do and do not resolve, maps each schema to current ProveIt Lean APIs, and supplies dependency-ordered formalization targets.  The archive also contains exact Python reference algorithms, fourteen passing tests, worked JSON certificates, a theorem ledger, a difficulty matrix, and source-oriented Lean files.  The result is intended as an integration blueprint: it does not claim that the new generic compiler has already been accepted by Lean's kernel.

\end{abstract}

\vspace{0.5em}
\noindent\textbf{Verification notice.}
This report distinguishes four statuses throughout: inherited Lean-kernel-checked results already in ProveIt, new human proofs, uncompiled Lean blueprints, and theorem-specific residual producer obligations.  The accompanying Python code is an exact semantic reference implementation and test suite; it is not an object-language proof checker.  No new Lean-kernel theorem is claimed by this archive.

\tableofcontents
\newpage

\section{Repository baseline and the verification boundary}
\label{sec:baseline}

\subsection{The three repository layers}

The present repository should be read as three logically distinct layers.

\begin{enumerate}[label=\textbf{L\arabic*.},leftmargin=3.2em]
\item \textbf{A checked object calculus.}  \texttt{ConciseIPRA} contains the ordinary first-order presentation \(\Prf\;\PRAAx\), together with the native scalar presentations and their bi-interpretations.  The ordinary presentation is the right portability boundary: formulas, local contexts, substitution, weakening, equality, quantifiers, and quantifier-free induction are explicit.

\item \textbf{A checked arithmetic library.}  \texttt{IPRADerivations} contains actual values of type
\[
  ([] : \mathsf{List}(\mathsf{Fm}\;\mathsf{PTm}\;0))
  \;\vdash[\PRAAx]\;\varphi.
\]
Its reusable core already includes semiring laws, order, monus, pairing, simultaneous recursion, bounded sums and products, bounded least search, division, gcd, Bezout, finite linear algebra, and several forms of descent.

\item \textbf{Research packets.}  The proof-packet collection contains broad human reductions and Lean plans for many mathematical topics.  It is a source of definitions, quantitative reformulations, examples, and obstacle analyses.  It is not, merely by being in the repository, a checked extension of the object theory.  The repository README explicitly records that the packets were filed as delivered and that many Lean files are uncompiled specifications or outlines.
\end{enumerate}

The framework developed here uses L1 and L2 as trusted implementation dependencies and treats L3 as prior research to be consolidated, corrected, or specialized.  In particular, a host-level theorem about natural numbers, a Python test, and a theorem whose result type is literally \(\Prf\;\PRAAx\;\Gamma\;\varphi\) are different deliverables.

\subsection{Canonical endpoint}

For new portable work, the canonical endpoint is
\begin{equation}
  \Gamma\;\vdash[\PRAAx]\;\varphi,
  \label{eq:canonical-endpoint}
\end{equation}
where \(\Gamma\) is an explicit finite list of object-language formulas and \(\varphi\) is an explicit object-language formula over the ordinary \texttt{PTm}/\texttt{Fm} syntax.  The scalar calculi are valuable for proof-checker minimization, but theorem development is easier to share through the ordinary calculus.  Foundation-side consumers may receive a translated corollary through the bridge after the ordinary endpoint exists.

This choice has four consequences.

\begin{enumerate}[leftmargin=2.2em]
\item Every mathematical operation used by the proof is represented by a fixed primitive-recursive symbol or by a fixed term assembled from such symbols.
\item Every induction invariant is quantifier-free in the object language.  Bounded quantifiers are first compiled to Boolean folds.
\item Every unproved quantitative input appears in \(\Gamma\); it is not hidden in a Lean structure, a Python value, or an informal phrase such as ``choose the constant sufficiently large.''
\item Reuse happens at arbitrary terms and under arbitrary local contexts whenever possible.  Closed theorems are useful, but open lemmas are the compositional currency.
\end{enumerate}

\subsection{Checked kernel operations already sufficient for composition}

The following existing rules are enough to support the structural part of the proposed compiler.

\begin{description}[style=nextline,leftmargin=2.2em]
\item[Substitution and weakening.] If \(\Gamma\vdash\varphi\), then simultaneous substitution of object terms gives the substituted sequent; the context may also be weakened.

\item[Instantiation of closed derivations.]  The checked theorem \texttt{LeanProofs.IPRA.inst} specializes a closed derivation with free parameters at arbitrary terms under an arbitrary context.

\item[Quantifier-free induction.]  The checked theorem \texttt{LeanProofs.IPRA.ind} packages the \iPRA{} induction scheme for an open quantifier-free formula while retaining the other variables as parameters.

\item[Simultaneous cut.]  The checked theorem \texttt{Prf.trans} states that if \(\Gamma\vdash\varphi\) and every member of \(\Gamma\) is derivable from \(\Delta\), then \(\Delta\vdash\varphi\).  This is exactly the operation needed to close a residualized compiler endpoint.

\item[Bounded least search.]  \texttt{Core.Search} defines a primitive-recursive least-search term and proves its bound, success, and minimality contracts in the object theory.
\end{description}

Thus the largest missing component is not a stronger logic.  It is a library of proof-producing adapters that normalize finite mathematics to these existing operations.

\subsection{What this report claims}

The mathematical schemas in this report are proved at the human level with enough detail to guide a line-by-line formalization.  The archive's \texttt{lean/} directory then states implementation targets against the current API.  Except for a tiny wrapper file around already checked declarations, those files are explicitly blueprints and are not claimed to elaborate.  The Python companion demonstrates exact finite algorithms and independently checkable certificates; it supplies no additional axiom and no object proof.

This status discipline is part of the framework, not an administrative afterthought.  A compiler for weak arithmetic is useful only if one can see, for every theorem, which pieces have reached the kernel and which pieces remain mathematical producers.

\section{The FINER normal form}
\label{sec:design}

\subsection{The slogan: query first, evidence second, object proof last}

Many ordinary mathematical statements quantify over completed objects: real functions, infinite sequences, compact spaces, measures, branches of algebraic functions, or all stages of a construction.  Such a statement is usually not the right first object for \iPRA.  The correct first object is a \emph{finite query}.  A query asks for exactly the finite observation needed by the surrounding theorem: a precision, a height, a time horizon, a degree, a mesh, a finite diagram, or a bounded list of tests.

For a query code \(q\), FINER constructs five theorem-specific objects:
\begin{equation}
  M(q),\qquad \CertCheck(q,c),\qquad \Obs(q,c),\qquad
  \Resid_1(q),\ldots,\Resid_r(q),
  \label{eq:finer-data}
\end{equation}
where \(M\), \(\CertCheck\), and \(\Obs\) are primitive recursive and \(c<M(q)\) is a single master certificate.  The central object-language theorem has the shape
\begin{equation}
 \Resid_1(q),\ldots,\Resid_r(q),\Gamma(q)
 \;\vdash[\PRAAx]\;
 \Phi\bigl(q,\Obs(q,\Sel(q))\bigr),
 \label{eq:finer-endpoint}
\end{equation}
with
\begin{equation}
 \Sel(q)=\mu c<M(q)\,[\CertCheck(q,c)=1].
 \label{eq:finer-selector}
\end{equation}
The least-search operation in \eqref{eq:finer-selector} is not an informal choice and is not a new axiom.  It is a fixed primitive-recursive program with an existing object-theory correctness proof.

\subsection{The six node classes}

A specialized finite proof graph is normalized to six node classes.  These classes are deliberately less expressive than a programming language: their purpose is to expose proof obligations.

\begin{enumerate}[label=\textbf{N\arabic*.},leftmargin=3.2em]
\item \textbf{Exact computation.}  A fixed primitive-recursive transducer computes a term or a finite array.  Its local certificate consists of defining equations.

\item \textbf{Boolean decision.}  A fixed bounded formula is reflected by a bit.  The local certificate proves bit normalization and the two reflection directions.

\item \textbf{Bounded search.}  A finite search returns the least accepted code or the bound on failure.  Its local certificate is obtained from the checked search library.

\item \textbf{Certified iteration.}  A clocked loop is represented by primitive recursion.  Its safety or correctness invariant is quantifier-free, so one application of \iPRA{} induction reflects the run.

\item \textbf{Finite relational interaction.}  A producer may supply several locally valid answers.  Backward viability determines whether a successful bounded continuation exists; forward least search then produces one transcript.  This avoids a global selector.

\item \textbf{Observation.}  Internal certificates are erased only through an exact invariant, an explicit quotient congruence, a metric diameter bound, or a guard such as a positive separation margin.
\end{enumerate}

A source proof may contain cyclic precision dependencies.  FINER resolves these before execution by the finite monotone demand closure of \cref{sec:demand}.  Once the finite demands are fixed, the execution and justification graph is unrolled to a finite schedule.

\subsection{No universal evaluator}

The graph in FINER is \emph{metalevel finite data}.  Specialization generates a fixed \texttt{PR} syntax tree for each transducer and checker.  The object theory does not contain a total primitive-recursive interpreter for arbitrary primitive-recursive programs.

\begin{barrier}[Diagonal obstruction]
There is no primitive-recursive function \(U(e,x)\) that enumerates every unary primitive-recursive function as \(x\mapsto U(e,x)\).
\end{barrier}

\begin{proof}
If such a \(U\) existed, then \(d(x)=U(x,x)+1\) would itself be primitive recursive, hence \(d(x)=U(e,x)\) for some \(e\).  Evaluating at \(e\) gives
\[
  U(e,e)+1=d(e)=U(e,e),
\]
contradicting successor injectivity and the fact that no successor is zero.  The contradiction can be rendered in very weak arithmetic once the alleged universal evaluator and its enumeration property are fixed.
\end{proof}

Consequently, a generic compiler theorem is a metatheorem about generating object syntax and proofs from a fixed graph.  Inside Lean, it may be represented by an inductive metalevel graph whose compiler returns \texttt{PR}, \texttt{PTm}, \texttt{Fm}, and \texttt{Prf} values.  It must not be confused with an object-language universal machine.

\subsection{Two horizons}

Every nontrivial finite node has two resource measures.

\begin{definition}[Computation and justification horizons]
The \emph{computation horizon} bounds the data actually inspected or generated.  The \emph{justification horizon} bounds the evidence needed to prove that the computation has the requested property.
\end{definition}

The two can differ.  An approximate root computation may inspect \(k\) bisection bits but require a derivative or separation estimate at a finer precision.  A finite cover computation may output one index but require a mesh and a modulus that justify coverage.  Keeping the two horizons explicit prevents a common circularity: increasing precision to justify a step can itself require stronger justification.  The demand-closure theorem computes the least simultaneous finite solution whenever all dependencies are monotone and capped.

\subsection{Status as part of the type}

FINER recommends that each node be tagged with one of three proof statuses.

\begin{description}[leftmargin=3em,style=nextline]
\item[Derived.] The local sequent is already a checked \(\Prf\;\PRAAx\) value.
\item[Residual.] The local fact is retained as an explicit hypothesis of the final sequent.
\item[External only.] A host-language theorem or executable test is available, but no object-language import has been built.  Such evidence may guide development but cannot silently populate either of the first two classes.
\end{description}

The compiler may compose Derived nodes and carry Residual nodes.  It must reject an External-only node as an object proof dependency until an adapter is provided.  This simple rule prevents most accidental overclaims found in broad formalization blueprints.

\section{Finite data, coding, and local contracts}
\label{sec:finite-data}

\subsection{One numerical sort is enough}

All finite data used by the framework are natural numbers.  Tuples, signed integers, rationals, arrays, finite sets, graphs, trees, matrices, formulas of a fixed grammar, and transcripts are represented by primitive-recursive codes.  There is no requirement that all applications share one global encoding.  A theorem-specific encoding is often smaller and easier to verify.

The checked repository already supplies a pairing function
\[
  \langle a,b\rangle=(a+b)^2+a
\]
with primitive-recursive projections and object proofs of the projection equations.  Repeated pairing is adequate for fixed arity.  For arrays whose length is a query parameter, a mixed-radix or bounded-sequence representation is more convenient.

\subsection{Mixed-radix arrays}

Let \(r_0,\ldots,r_{m-1}\geq 1\) be radices and let \(0\leq d_i<r_i\).  Define
\begin{equation}
  \code(d_0,\ldots,d_{m-1})
  =\sum_{i<m} d_i\prod_{j<i}r_j.
  \label{eq:mixed-code}
\end{equation}
The total range is
\[
  R=\prod_{i<m}r_i.
\]
The \(i\)-th digit is obtained by quotienting by \(\prod_{j<i}r_j\) and taking remainder modulo \(r_i\).  All operations are primitive recursive.  The checked division-with-remainder library supplies the arithmetic needed for the object proof.

\begin{lemma}[Mixed-radix correctness]
For every bounded digit vector, \(\decode_i(\code(\vec d))=d_i\).  Conversely, if \(c<R\), then re-encoding its decoded digits gives \(c\).
\end{lemma}

\begin{proof}
For the first direction, split the sum in \eqref{eq:mixed-code} below, at, and above \(i\).  Terms below \(i\) form a number smaller than the place value \(P_i=\prod_{j<i}r_j\).  The term at \(i\) is \(d_iP_i\).  Every term above \(i\) is divisible by \(P_ir_i\).  Division by \(P_i\), followed by remainder modulo \(r_i\), therefore leaves exactly \(d_i\).

For the converse, repeatedly apply the quotient-remainder equation
\[
  c=r_0(c\mathbin{\mathrm{quo}}r_0)+(c\mathbin{\mathrm{rem}}r_0).
\]
After \(m\) steps the remaining quotient is zero because \(c<R\).  Substitution yields the expansion \eqref{eq:mixed-code}.  In \iPRA, the variable-length statement is proved by quantifier-free induction on \(m\), with the current quotient and place value carried in the state.
\end{proof}

Mixed radix makes a master certificate literal: each node receives one bounded digit block, and all blocks form one number below a primitive-recursive product.  It also makes certificate minimality meaningful, because the search order is fixed.

\subsection{Finite predicates and characteristic bits}

A predicate used inside a loop is represented by a normalized bit \(b\) satisfying
\begin{equation}
  b=0\;\vee\;b=1.
  \label{eq:bit-normal}
\end{equation}
The repository's equality bit \texttt{eqb} and order development provide the atomic cases.  Boolean connectives are fixed primitive-recursive operations, for example
\[
\begin{aligned}
  \mathsf{not}(b)&=1\mathbin{\dot{-}} b,\\
  \mathsf{and}(b,c)&=bc,\\
  \mathsf{or}(b,c)&=\mathsf{sg}(b+c).
\end{aligned}
\]
Because inputs are normalized, their truth tables are derivable by finite case distinction in the intuitionistic object theory.

\subsection{Local proof contracts}

\begin{definition}[Finite contract]
A finite contract for queries \(q\) and outputs \(y\) consists of:
\begin{enumerate}[label=(\roman*),leftmargin=2.3em]
\item a certificate bound \(M(q)\);
\item a Boolean checker \(C(q,c)\);
\item an output decoder \(O(q,c)\);
\item residual formulas \(B_1(q),\ldots,B_r(q)\);
\item an object derivation
\[
  C(q,c)=1,\;B_1(q),\ldots,B_r(q),\Gamma(q)
  \;\vdash[\PRAAx]\;\Phi(q,O(q,c)).
\]
\end{enumerate}
All displayed programs are fixed primitive-recursive symbols after theorem specialization.
\end{definition}

A contract deliberately separates \emph{soundness} from \emph{inhabitation}.  The checker may be perfectly sound even when no accepted certificate exists.  Inhabitation is either derived by a producer theorem or retained as the residual
\begin{equation}
  \exists c<M(q)\;[C(q,c)=1].
  \label{eq:certificate-exists}
\end{equation}
This distinction is essential in Diophantine geometry, effective analysis, and compactness arguments: a finite verifier does not create a height bound, modulus, or cover.

\subsection{Composition of contracts}

Suppose contract \(A\) produces \(y\) from \(q\), and contract \(B\) produces \(z\) from \((q,y)\).  A composite master certificate is the pair \(\langle c_A,c_B\rangle\), or the corresponding mixed-radix block.  Its checker is
\[
 C_A(q,c_A)\wedge C_B(q,O_A(q,c_A),c_B),
\]
and its output is \(O_B(q,O_A(q,c_A),c_B)\).  The proof is a substitution of the first soundness theorem into the second, followed by simultaneous cut.  Residual lists concatenate, with duplicates optionally removed by a metalevel normalization.

No categorical machinery is needed for this basic operation, though finite contracts do form a useful category at the metalevel.  The crucial point is that composition preserves the object-theory endpoint and exposes, rather than hides, the union of residual producers.

\section{Fixed bounded formulas and Boolean reflection}
\label{sec:boolean}

The induction scheme of \iPRA{} applies only to quantifier-free formulas.  This is not an obstacle to finite mathematics, provided bounded logical content is compiled to primitive-recursive bits before induction.  The word \emph{fixed} is important: compilation is structural recursion in Lean on a theorem-specific syntax tree, not evaluation of an arbitrary object-language program.

\subsection{A fixed bounded grammar}

For a fixed number of free variables, consider formulas generated by
\[
  A ::= s=t\mid s<t\mid A\wedge A\mid A\vee A\mid A\to A
       \mid \exists i<N\,A(i)\mid \forall i<N\,A(i),
\]
where all terms are fixed primitive-recursive terms.  Negation is \(A\to\bot\).  Bounds may depend on the outer free variables but not on the variable they bind.

\begin{theorem}[Boolean reflection]
\label{thm:boolean-reflection}
For every fixed bounded formula \(A(\vec x)\), one can effectively construct a primitive-recursive term \(\chi_A(\vec x)\) and object derivations, under every context \(\Gamma\), of
\begin{align}
 &\chi_A=0\;\vee\;\chi_A=1, \label{eq:reflect-bit}\\
 &\chi_A=1\;\leftrightarrow\; A, \label{eq:reflect-one}\\
 &\chi_A=0\;\leftrightarrow\; \neg A. \label{eq:reflect-zero}
\end{align}
The construction is uniform at the metalevel in the syntax tree of \(A\).  It introduces no new nonlogical axiom.
\end{theorem}

\begin{proof}
We proceed by structural recursion on \(A\).

\paragraph{Equality.}
Take \(\chi_{s=t}=\mathsf{eqb}(s,t)\).  The checked equality-bit development proves \(\mathsf{eqb}(s,t)=1\) from \(s=t\), and \(\mathsf{eqb}(s,t)=0\) from \(s\neq t\).  The checked totality of equality gives \(s=t\vee s\neq t\), so the converse directions and bit normalization follow by disjunction elimination.

\paragraph{Order.}
Use the characteristic bit for the repository's chosen order definition.  Order is expressed through monus and equality, so its normalization and reflection reduce to the preceding case plus the checked order theorems.

\paragraph{Conjunction.}
Let \(b=\chi_B\), \(c=\chi_C\), and set \(\chi_{B\wedge C}=bc\).  Four finite cases from \eqref{eq:reflect-bit} for \(b,c\) establish normalization.  If \(bc=1\), neither factor can be zero; the Boolean cases force \(b=c=1\), and the induction hypotheses give \(B\wedge C\).  Conversely, \(B\wedge C\) yields \(b=c=1\), hence \(bc=1\).  The zero case is similar and intuitionistically valid because both component formulas are decidable.

\paragraph{Disjunction.}
Set \(\chi_{B\vee C}=\mathsf{sg}(b+c)\), where \(\mathsf{sg}(0)=0\) and \(\mathsf{sg}(S n)=1\).  If one input bit is one, the sum is a successor and the output bit is one.  If the output is one, the sum is nonzero; finite Boolean case analysis excludes \(b=c=0\), so one component formula holds.  The zero direction follows from \(b+c=0\), which implies both summands are zero.

\paragraph{Implication.}
Define
\[
  \chi_{B\to C}=\mathsf{or}(\mathsf{not}(b),c).
\]
Classically one writes \(B\to C\equiv\neg B\vee C\).  Here that equivalence is justified intuitionistically because \(B\) is decidable by the induction hypothesis: split on \(B\vee\neg B\).  In the first branch an implication gives \(C\); in the second branch \(\neg B\) is already available.  The reverse direction is intuitionistically immediate.  The Boolean truth table then gives \eqref{eq:reflect-one}--\eqref{eq:reflect-zero}.

\paragraph{Bounded existential.}
Assume \(B(i,\vec x)\) has reflected bit \(b(i,\vec x)\).  Define a finite sum
\[
  S(N,\vec x)=\sum_{i<N}b(i,\vec x)
\]
and put \(\chi_{\exists i<N B}=\mathsf{sg}(S)\).  The defining equations of the checked finite-sum symbol are available at arbitrary terms.  Quantifier-free induction on \(N\) proves
\begin{align*}
  S(N,\vec x)=0 &\;\to\; \forall i<N\;b(i,\vec x)=0,\\
  \bigl(\exists i<N\;b(i,\vec x)=1\bigr) &\;\to\; S(N,\vec x)\neq 0.
\end{align*}
For witness extraction from \(S\neq 0\), one may either use the same induction or, more modularly, invoke bounded least search on the zero-body \(1\dot{-} b(i,\vec x)\).  The search theorem returns the least index with \(b=1\).  Reflection for \(B\) converts between the bit and formula witnesses.

\paragraph{Bounded universal.}
There are two convenient definitions.  One may use the product
\[
  P(N,\vec x)=\prod_{i<N}b(i,\vec x)
\]
after proving that all factors are bits, or define the universal bit as the negation of the existential failure bit:
\[
  \chi_{\forall i<N B}
  =\mathsf{not}\bigl(\chi_{\exists i<N\neg B}\bigr).
\]
The second choice reuses the existential development.  Because every instance \(B(i,\vec x)\) is decidable, intuitionistic bounded De Morgan reasoning proves
\[
 \neg\exists i<N\neg B(i)\;\leftrightarrow\;\forall i<N B(i).
\]
The proof is a quantifier-free induction on \(N\), or a bounded least-counterexample argument.

Each constructor preserves substitution because the generated terms and formulas are built from syntax whose substitution laws are already proved.  Thus the result is usable under arbitrary local hypotheses and at arbitrary terms.
\end{proof}

\subsection{Least witnesses and counterexamples}

Reflection is more useful when it comes with extraction.

\begin{corollary}[Least bounded witness]
\label{cor:least-bounded-witness}
For every fixed reflected body \(B(i,\vec x)\), there is a primitive-recursive term \(\mu_B(N,\vec x)\) such that \iPRA{} derives:
\begin{enumerate}[label=(\roman*),leftmargin=2.3em]
\item \(\mu_B\leq N\);
\item if \(\exists i<N\,B(i,\vec x)\), then \(\mu_B<N\), \(B(\mu_B,\vec x)\), and \(\mu_B\leq i\) for every supplied witness \(i<N\);
\item if \(j<\mu_B\), then \(\neg B(j,\vec x)\).
\end{enumerate}
\end{corollary}

\begin{proof}
Apply the checked \texttt{Search.mu} to the term
\[
  f(i,\vec x)=1\dot{-}\chi_B(i,\vec x).
\]
For normalized bits, \(f=0\) is equivalent to \(\chi_B=1\), hence to \(B\).  The statements are direct instances of \texttt{search\_spec}, \texttt{search\_found}, and \texttt{search\_minimal}, followed by the reflection equivalence.
\end{proof}

The dual least-counterexample term is obtained by searching for \(\chi_B=0\).  This is the basis of finite failure certificates and universal strategy checking.

\subsection{Fixed alternation and finite strategies}

A fixed bounded quantifier prefix
\[
  Q_0x_0<B_0\;Q_1x_1<B_1\cdots Q_{m-1}x_{m-1}<B_{m-1}\;A(\vec x)
\]
can be evaluated backwards as a finite game.  Existential levels take a finite disjunction; universal levels take a finite conjunction.  A success certificate stores one move at existential levels and all required child certificates at universal levels.  A failure certificate reverses the roles.  Since \(m\) and the source formula are fixed, the resulting verifier is primitive recursive and its soundness is proved by induction on the fixed syntax tree plus quantifier-free induction over variable bounds.

This recovers the useful part of finite classical reasoning: every bounded position has a decidable winning bit and one of two finite strategy certificates.  It does \emph{not} recover an unbounded choice function or a witness for an arbitrary classically proved existential statement.

\subsection{Lean implementation route}

A practical Lean implementation should define a small metalevel datatype \texttt{BFormula n} with constructors for atoms, connectives, and bounded quantifiers.  Compilation returns a structure containing a \texttt{PTm n}, an \texttt{Fm PTm n}, and the five derivations in \cref{thm:boolean-reflection}.  The constructors should call existing arithmetic lemmas rather than reprove Boolean algebra.  In particular:

\begin{itemize}[leftmargin=2em]
\item equality uses \texttt{eqb\_of\_eq}, \texttt{eqb\_of\_ne}, and \texttt{eq\_or\_ne};
\item bounded witnesses use \texttt{Search.search\_found};
\item folds use \texttt{sumF}/\texttt{prodF} and their equations;
\item the only induction is \texttt{LeanProofs.IPRA.ind} on a quantifier-free accumulator invariant.
\end{itemize}

The generated \texttt{PR} symbol is specific to the fixed \texttt{BFormula}.  There is no object-level syntax code or evaluator unless a particular application independently requires one.

\section{Master certificates and checked selection}
\label{sec:certificates}

\subsection{Why one certificate matters}

A large finitary proof often introduces many local witnesses: a trace, a finite net, branch choices, error allocations, local algebraic identities, a spanning tree, and a list of bounded solutions.  Treating these as unrelated existential quantifiers invites hidden choice and makes the final theorem difficult to audit.  FINER instead assigns every local field a finite radix and packs all fields into one number \(c<M(q)\).

The checker \(\CertCheck(q,c)\) performs four kinds of tests:
\begin{enumerate}[label=(\roman*),leftmargin=2.3em]
\item digit and array range checks;
\item local defining-equation and arithmetic checks;
\item edge checks linking dependent node certificates;
\item terminal checks connecting the decoded observation to the target formula.
\end{enumerate}
The checker is normalized to a Boolean bit and reflected by \cref{thm:boolean-reflection}.  Therefore the semantic-looking statement ``the certificate passes'' is an object-language equation \(\CertCheck(q,c)=1\).

\subsection{Least accepted code}

Let
\[
  f_q(c)=1\dot{-}\CertCheck(q,c).
\]
The checked search module returns
\[
  \Sel(q)=\mu c<M(q)\,[f_q(c)=0].
\]
For a normalized checker bit, \(f_q(c)=0\) iff \(\CertCheck(q,c)=1\).

\begin{theorem}[Selected-certificate theorem]
\label{thm:selected-certificate}
Assume the object theory derives
\begin{equation}
  \exists c<M(q)\;[\CertCheck(q,c)=1].
  \label{eq:accepted-exists}
\end{equation}
Then it derives
\begin{align}
  &\Sel(q)<M(q),\label{eq:selected-bound}\\
  &\CertCheck(q,\Sel(q))=1,\label{eq:selected-accept}\\
  &c<\Sel(q)\;\to\;\CertCheck(q,c)=0.\label{eq:selected-minimal}
\end{align}
\end{theorem}

\begin{proof}
From \eqref{eq:accepted-exists}, existential elimination supplies a witness \(c\) and proofs of \(c<M(q)\) and \(\CertCheck(q,c)=1\).  Bit normalization converts the latter to \(f_q(c)=0\).  Apply \texttt{Search.search\_found}; it gives that the search result is below the bound, is a zero of \(f_q\), and is no larger than the supplied witness.  Convert \(f_q(\Sel(q))=0\) back to \(\CertCheck(q,\Sel(q))=1\).  Minimality is \texttt{Search.search\_minimal} with the same conversion.
\end{proof}

The leastness is not mathematically necessary for soundness, but it has three practical benefits.  It makes outputs deterministic, gives a canonical regression fixture, and permits exact comparison of independently generated certificates.

\subsection{Success and failure evidence}

When a bounded problem is decidable, the framework can expose both outcomes.

\begin{definition}[Dual certificate interface]
A dual interface has a success checker \(C^+\), a failure checker \(C^-\), and object proofs that:
\begin{align*}
 C^+(q,c)=1 &\to \Phi(q),\\
 C^-(q,d)=1 &\to \neg\Phi(q),\\
 \bigl(\exists c<M^+(q)\,C^+(q,c)=1\bigr)
 &\vee
 \bigl(\exists d<M^-(q)\,C^-(q,d)=1\bigr).
\end{align*}
\end{definition}

For a fixed bounded formula, the Boolean reflection construction supplies such an interface.  For a bounded existential, success is a witness and failure is a finite table of rejected candidates, or simply the reflected universal statement that every candidate fails.  For a finite game, success and failure are dual strategy trees.  The important property is not that every theorem has a short negative certificate, but that bounded classical case analysis becomes finite evidence on both branches.

\subsection{Certificate soundness is compositional}

Suppose a graph has nodes \(v_0,\ldots,v_{m-1}\) in a dependency order after demand closure and finite unrolling.  Let \(A_i(q,c)\) be the assertion that node \(i\)'s decoded block is valid.  The master checker is a Boolean conjunction of their reflected bits.  Reflection yields each \(A_i\) from master acceptance.

For every node, provide a local sequent
\begin{equation}
  A_i,\;\{P_j:j\in\operatorname{pred}(i)\},\;\Gamma_i
  \;\vdash[\PRAAx]\;P_i,
  \label{eq:local-node-sequent}
\end{equation}
where \(P_j\) are predecessor postconditions and \(\Gamma_i\) are residual assumptions used locally.  Repeated simultaneous cut along the finite order yields the terminal postcondition.  At no point does the compiler need to inspect the semantic meaning of \(P_i\); it only substitutes terms and composes derivations.

This is the proof-theoretic reason for using one master certificate.  Its accepted equation is a single reusable major premise from which all node facts are projected.

\subsection{Bounds and enormous finite spaces}

The certificate space may be astronomically large.  That does not invalidate the object theorem.  Primitive recursive arithmetic concerns total finite computations, not feasible ones.  Nevertheless, practical formalization benefits from factorization:

\begin{itemize}[leftmargin=2em]
\item deterministic fields should be recomputed rather than stored;
\item independent blocks should use local searches whose results are then packed;
\item bounds should be stated symbolically rather than expanded as numerals;
\item proof checking should use reflected local equations, not execute the entire search in Lean's kernel.
\end{itemize}

The role of least search is logical witness extraction.  An external producer may directly supply a compact certificate for testing and replay; the object theorem remains valid because it proves the least accepted certificate exists whenever any accepted certificate exists.

\section{Trace reflection with one quantifier-free induction}
\label{sec:trace}

Many proofs in the packets use a finite algorithm and then argue informally that each step preserves an invariant.  FINER isolates a generic object-theory theorem for this pattern.

\subsection{Deterministic state recursion}

Fix primitive-recursive terms
\[
  s_0(q,c),\qquad \mathsf{step}(t,s,q,c),
\]
and define the state sequence by primitive recursion:
\begin{align*}
 S(0,q,c)&=s_0(q,c),\\
 S(t+1,q,c)&=\mathsf{step}(t,S(t,q,c),q,c).
\end{align*}
Let \(I(t,s,q,c)\) be a fixed bounded formula.  Compile it to a bit \(\chi_I\) and use the quantifier-free invariant
\begin{equation}
  \chi_I(t,S(t,q,c),q,c)=1.
  \label{eq:qf-trace-invariant}
\end{equation}

\begin{theorem}[Deterministic trace reflection]
\label{thm:trace-reflection}
Suppose \iPRA{} derives, at arbitrary terms and under a context \(\Gamma\),
\begin{align}
 &\chi_I(0,s_0(q,c),q,c)=1,\label{eq:trace-base}\\
 &\chi_I(t,s,q,c)=1\to
   \chi_I(t+1,\mathsf{step}(t,s,q,c),q,c)=1.
   \label{eq:trace-step}
\end{align}
Then it derives
\begin{equation}
  \chi_I(T,S(T,q,c),q,c)=1
  \label{eq:trace-end}
\end{equation}
for every term \(T\).  If a terminal lemma derives \(\Phi(q,c)\) from \eqref{eq:trace-end}, then \(\Gamma\vdash\Phi(q,c)\).
\end{theorem}

\begin{proof}
The formula in \eqref{eq:qf-trace-invariant} is an equation of primitive-recursive terms, hence quantifier-free.  Apply the checked open induction theorem \texttt{ind}.  The base case reduces by the zero equation of the state recursor to \eqref{eq:trace-base}.  In the successor case, the induction hypothesis is \(\chi_I(t,S(t,q,c),q,c)=1\).  Instantiate \eqref{eq:trace-step} with \(s=S(t,q,c)\), then rewrite its conclusion using the successor equation for \(S\).  Universal closure and instantiation give the result at arbitrary \(T\).  The terminal lemma is composed by implication elimination or simultaneous cut.
\end{proof}

The theorem is elementary, but making it a reusable object derivation removes a large amount of theorem-specific boilerplate.

\subsection{Checked transcripts}

A nondeterministic or externally generated run is represented by a packed transcript \(\tau\).  Let \(\mathsf{state}(\tau,t)\) decode its \(t\)-th state.  The transcript checker verifies
\begin{align*}
  \mathsf{state}(\tau,0)&=s_0,\\
  R(t,\mathsf{state}(\tau,t),\mathsf{state}(\tau,t+1))&=1
  \quad(t<T),\\
  A(\mathsf{state}(\tau,T))&=1.
\end{align*}
All bounded universal content is reflected to a fold bit.  Soundness can be proved in either of two ways.

\begin{enumerate}[label=(\alph*),leftmargin=2.3em]
\item Use Boolean reflection to extract each edge equation on demand.  This is simplest when only finitely many fixed positions are used by a terminal proof.
\item Prove by quantifier-free induction that every prefix up to \(t\) is valid.  The state consists of \(t\), the decoded current state, and an accumulated validity bit.  This is preferable for variable horizons.
\end{enumerate}

The transcript need not be trusted.  It is merely a finite input whose accepted bit implies the edge formulas in the object theory.

\subsection{Normalized trace kernels}

To minimize future Lean work, a node-specific trace proof should be normalized to four derivations:
\begin{enumerate}[label=(\roman*),leftmargin=2.3em]
\item the state recursor's zero equation;
\item its successor equation;
\item a base invariant derivation;
\item a step invariant derivation under the invariant hypothesis.
\end{enumerate}
A generic theorem then produces the terminal invariant.  This design is preferable to a giant tactic that interprets arbitrary source code.  It keeps the trusted object proof small and lets theorem-specific normalizers remain ordinary Lean programs whose output is checked by the existing calculus.

\subsection{Changing parameters}

Some loops refine both a state and a demand parameter.  The correct state is their pair.  If
\[
  (d_{t+1},s_{t+1})=F(t,d_t,s_t),
\]
then pair the values and apply \cref{thm:trace-reflection}.  An invariant may mention projections, resource bounds, and an error budget, provided the final compiled formula is quantifier-free.  Simultaneous recursion from the checked pairing library is particularly useful when separate readable symbols for the components are desired.

\subsection{Nested loops}

Nested finite loops do not require nested logical induction if their combined machine is flattened.  Encode the program counter, local indices, and accumulators in one state; one outer clock bounds all transitions.  Alternatively, prove the inner loop once as a local contract and use it as the step lemma of the outer loop.  The second style usually yields smaller terms and more reusable lemmas.  In both cases, the logical requirement remains one quantifier-free induction per certified iterator, not induction over a formula containing bounded quantifiers.

\section{Finite monotone demand closure}
\label{sec:demand}

Precision propagation is often presented as a reverse pass over an acyclic graph.  That is insufficient when two modules demand finer data from one another.  The apparent cycle need not be fatal: if demands are bounded, monotone, and combined by maximum, the least simultaneous solution is obtained by finite iteration.

\subsection{Finite product of chains}

Fix a dimension \(m\) and capacities \(C_0,\ldots,C_{m-1}\).  Let
\[
  D=\prod_{i<m}\{0,1,\ldots,C_i\}
\]
with componentwise order.  Write \(x\vee y\) for componentwise maximum.  Let \(F:D\to D\) be monotone:
\begin{equation}
  x\leq y\quad\Longrightarrow\quad F(x)\leq F(y).
  \label{eq:demand-monotone}
\end{equation}
Define
\begin{equation}
  x^0=\vec 0,
  \qquad
  x^{t+1}=x^t\vee F(x^t).
  \label{eq:demand-iteration}
\end{equation}
Let
\begin{equation}
  H=\sum_{i<m}C_i,
  \qquad
  \rank(x)=\sum_{i<m}x_i.
  \label{eq:demand-rank}
\end{equation}

\begin{lemma}[Inflation and bounded rank]
\label{lem:demand-inflation}
For every \(t\), \(x^t\leq x^{t+1}\) and \(\rank(x^t)\leq H\).  If \(x^t\neq x^{t+1}\), then
\[
  \rank(x^{t+1})\geq \rank(x^t)+1.
\]
\end{lemma}

\begin{proof}
Inflation is immediate from the definition of componentwise maximum.  Every coordinate of every iterate lies between zero and its capacity because both arguments of the join do.  Summing gives the rank bound.  If the vectors differ, inflation implies that some coordinate strictly increases.  Natural-number strict increase means an increase by at least one; all other coordinates do not decrease.  Summing the coordinate inequalities gives the strict rank inequality.
\end{proof}

\begin{theorem}[Finite stabilization]
\label{thm:demand-stabilization}
The iteration is stable at time \(H\):
\begin{equation}
  x^{H+1}=x^H.
  \label{eq:demand-stable}
\end{equation}
Consequently \(F(x^H)\leq x^H\).
\end{theorem}

\begin{proof}
Assume for contradiction that \(x^{H+1}\neq x^H\).  We first show that every earlier transition is also strict.  If \(x^{t+1}=x^t\) for some \(t\), then \(F(x^t)\leq x^t\).  Hence the same update repeats \(x^t\), so by induction all later iterates are equal.  This would contradict strictness at \(H\).

Thus the transitions \(0\to1,1\to2,\ldots,H\to H+1\) are all strict.  Repeated application of \cref{lem:demand-inflation}, starting from rank zero, gives
\[
  \rank(x^{H+1})\geq H+1.
\]
But every vector in \(D\) has rank at most \(H\), contradiction.  Therefore \eqref{eq:demand-stable} holds.  Expanding the update equation gives
\[
  x^H\vee F(x^H)=x^H,
\]
which is equivalent componentwise to \(F(x^H)\leq x^H\).
\end{proof}

The proof deliberately uses \(H+1\) transitions.  This avoids a subtle off-by-one ambiguity when the rank reaches its maximum exactly at time \(H\).

\begin{theorem}[Least prefixed point]
\label{thm:demand-least}
Let \(y\in D\) satisfy \(F(y)\leq y\).  Then \(x^H\leq y\).  Hence \(x^H\) is the least prefixed point of \(F\), and also the least fixed point of the inflationary map \(G(x)=x\vee F(x)\).
\end{theorem}

\begin{proof}
We prove \(x^t\leq y\) by induction on \(t\).  The base \(\vec0\leq y\) is immediate.  If \(x^t\leq y\), monotonicity gives \(F(x^t)\leq F(y)\leq y\).  Both \(x^t\) and \(F(x^t)\) are below \(y\), so their join \(x^{t+1}\) is below \(y\).  Instantiate at \(H\).  By \cref{thm:demand-stabilization}, \(x^H\) is itself prefixed, so it is least.
\end{proof}

\subsection{Why the theorem is available in iPRA}

For fixed \(m\), all vector operations are tuples of primitive-recursive terms.  For variable \(m\), vectors are packed arrays, and componentwise comparison and sums are bounded folds.  The following formulas can therefore be made quantifier-free after Boolean reflection:
\begin{itemize}[leftmargin=2em]
\item every coordinate lies below its capacity;
\item the current iterate is componentwise below the next;
\item a strict vector change raises rank;
\item an arbitrary prefixed point bounds the current iterate.
\end{itemize}
The trace \(x^t\) is a primitive-recursive iterator.  \cref{lem:demand-inflation} and the leastness induction use the checked \texttt{ind} rule.  The finite sum is supplied by \texttt{sumF} or a specialized fixed-arity term.  The final contradiction is arithmetic order reasoning.

A particularly formalization-friendly invariant is
\begin{equation}
  \mathsf{Bounded}(x^t)\;\wedge\;
  x^t\leq y\;\wedge\;
  \rank(x^t)\geq \mathsf{strictCount}(t),
  \label{eq:demand-invariant}
\end{equation}
where \(\mathsf{strictCount}\) counts transitions whose equality bit is zero.  One then proves \(\mathsf{strictCount}(H+1)\leq H\).  An even smaller proof avoids storing the count: assume the final step is strict, derive by backward persistence that all earlier steps are strict, and use a forward rank induction.  The first version is more uniform for generated proofs; the second is shorter for a fixed graph.

\subsection{Capped transformers and invalid demands}

Real demand functions may propose values larger than their declared capacities.  Define
\[
  \widehat F_i(x)=\min(F_i(x),C_i).
\]
The clipped transformer is primitive recursive and monotone.  Its least closure exists by the preceding theorem.  However, clipping can conceal an unsatisfied request.  FINER therefore adds an overflow bit
\[
  \mathsf{overflow}(x)=1
  \quad\Longleftrightarrow\quad
  \exists i<m\;F_i(x)>C_i.
\]
A successful contract requires overflow zero at the closure.  If the bit is one, the compiler reports a residual asking for larger capacities or rejects the specialization; it does not pretend the clipped demand is adequate.

\subsection{From cyclic graph to finite schedule}

Let every graph edge carry a monotone demand transformer.  The global \(F\) takes, for each node, the maximum of all demands received from its consumers and the node's intrinsic minimum.  The closure \(d^*=x^H\) simultaneously fixes all query precisions, horizons, mesh levels, and certificate bounds.  Every node can then be specialized at \(d^*\).  A computational cycle is handled separately by an explicit clocked iterator or relational transcript; the demand cycle itself has disappeared into a finite certificate.

\begin{corollary}[Minimal adequate specialization]
If \(d\) satisfies every declared demand constraint, then \(d^*\leq d\).  Thus FINER does not silently over-strengthen residual producer obligations beyond the least capped solution.
\end{corollary}

\begin{proof}
The demand constraints say exactly \(F(d)\leq d\).  Apply \cref{thm:demand-least}.
\end{proof}

This minimality is useful for analytical estimates: a producer only has to meet the precision actually forced by the finite consumer, not a globally chosen all-purpose precision.

\subsection{Executable certificate}

The accompanying reference implementation returns the entire trace, the closure, and the stabilization index.  A separate verifier recomputes every join step, checks the final prefix condition, and checks the rank bound.  A brute-force routine verifies leastness on small examples.  These tests illustrate the theorem but do not replace its object-language proof.

\section{Relational witnesses and backward viability}
\label{sec:viability}

Exact mathematical witnesses are often non-unique, discontinuous, or expensive to name.  A producer may nevertheless return a finite relation of locally valid answers.  A bounded consumer only needs a compatible finite run.  Backward viability converts that need into a deterministic finite calculation without assuming a global choice function.

\subsection{Layered relational systems}

Fix a horizon \(T\).  At time \(t\), states satisfy \(s<B_t(q)\).  Let
\[
  R_t(q,s,s')\in\{0,1\}
\]
be a reflected edge relation from layer \(t\) to layer \(t+1\), and let \(A(q,s)\) be the terminal acceptance bit.  Define viability backwards:
\begin{align}
 V_T(q,s)&=A(q,s),\label{eq:viable-terminal}\\
 V_t(q,s)&=\bigvee_{s'<B_{t+1}(q)}
   \bigl(R_t(q,s,s')\wedge V_{t+1}(q,s')\bigr).
 \label{eq:viable-backward}
\end{align}
The disjunction is a finite Boolean fold.  Hence the complete table \(V_t(s)\) is primitive recursive.  For variable \(T\), compute it by recursion on the reversed time index, packing the current layer table into one number.

\begin{theorem}[Viability correctness]
\label{thm:viability-correct}
For every state \(s<B_t(q)\), \iPRA{} proves
\begin{equation}
 V_t(q,s)=1
 \quad\leftrightarrow\quad
 \text{there is an accepted relational path from \(s\) of length \(T-t\)}.
 \label{eq:viability-path-equivalence}
\end{equation}
Here the right side is represented by one bounded path code and a fixed primitive-recursive path checker.
\end{theorem}

\begin{proof}
Induct on the remaining horizon \(h=T-t\).

For \(h=0\), a path consists only of its initial state, and acceptance is exactly \eqref{eq:viable-terminal}.

For the successor case, expand \eqref{eq:viable-backward}.  If \(V_t(q,s)=1\), Boolean reflection of the finite disjunction yields some \(s'<B_{t+1}\) such that \(R_t(q,s,s')=1\) and \(V_{t+1}(q,s')=1\).  The induction hypothesis supplies a path from \(s'\); prepend \(s\).  Conversely, any accepted path has a first successor \(s'\) satisfying the edge relation and an accepted tail.  The induction hypothesis gives \(V_{t+1}(q,s')=1\), so the corresponding disjunct makes \(V_t(q,s)=1\).

The induction predicate is not the displayed existential formula.  In the object proof, both sides are reflected bits: one is the recursively computed viability bit and the other is the path-checker bit.  Their equality is quantifier-free.
\end{proof}

\subsection{Least viable replay}

Define the chosen successor of a viable state to be the least \(s'<B_{t+1}\) such that
\[
  R_t(q,s,s')=1\quad\text{and}\quad V_{t+1}(q,s')=1.
\]
The body is reflected, so \cref{cor:least-bounded-witness} supplies the term.  Starting from an initial viable state \(s_0\), primitive recursion generates a canonical path.

\begin{theorem}[Least successful transcript]
\label{thm:least-path}
If \(V_0(q,s_0)=1\), the forward least-viable replay produces a path \(s_0,\ldots,s_T\) such that every edge relation holds and \(A(q,s_T)=1\).
\end{theorem}

\begin{proof}
Use quantifier-free induction on \(t\) with invariant
\[
  V_t(q,s_t)=1
  \quad\wedge\quad
  \mathsf{PrefixValid}(q,t,\tau)=1.
\]
after compiling prefix validity to a bit.  The base is the hypothesis and the empty-prefix equation.  In the step, viability correctness guarantees at least one viable successor; least search returns one and proves both the edge bit and next-layer viability.  Extend the prefix checker.  At \(T\), \eqref{eq:viable-terminal} converts viability to acceptance.
\end{proof}

This theorem is a finite replacement for choice.  It does not select a witness independently at every precision or for every input object.  It selects the least answer needed by one bounded interaction, after proving that this answer has a successful continuation.

\subsection{Failure certificates}

If \(V_0(q,s_0)=0\), the backward table is finite evidence that no accepted path exists.  A checker verifies the terminal table and every backward recurrence.  Soundness follows from the reverse direction of \cref{thm:viability-correct}.

This dual evidence is useful in decision procedures.  A successful search returns a path; a failed search returns a table that closes all possible continuations.  Neither branch relies on an external assertion that the search was exhaustive.

\subsection{Bounded alternation}

Relational viability generalizes to fixed bounded games.  At an existential node, the winning bit is a finite disjunction of winning children; at a universal node, it is a finite conjunction.  A winning existential strategy stores one selected child at existential levels and child strategies for every universal response.  The dual certificate reverses these roles.

\begin{proposition}[Fixed finite game reflection]
For every fixed finite game grammar with primitive-recursive move bounds and reflected terminal atoms, there are primitive-recursive winning and losing bits, bounded strategy encodings, and object proofs that an accepted strategy has the claimed outcome.
\end{proposition}

\begin{proof}
Structural induction on the fixed game tree constructor.  Terminal nodes use atomic reflection.  Existential and universal nodes use the bounded disjunction and conjunction cases of \cref{thm:boolean-reflection}; witness extraction and least counterexamples use bounded search.  Strategy soundness follows by the same induction.  If the horizon is a query variable, represent the backward evaluator as a clocked state machine and apply trace reflection.
\end{proof}

This proposition captures finite classical search, minimax arguments, bounded adversarial constructions, and finite back-and-forth comparisons.  It does not yield determinacy or strategies for unbounded games.

\subsection{Compatible paths versus independent witnesses}

A common invalid step is:
\begin{quote}
For every \(t\) there exists a valid witness \(w_t\); therefore choose witnesses that form a valid run.
\end{quote}
Pointwise existence does not imply compatibility.  Backward viability repairs the argument by testing continuations before choosing the current answer.  If the initial state is not viable, the desired compatible run simply does not exist within the declared finite bounds, even though every layer may be nonempty.

For outputs that do not require a compatible run, a stronger and cheaper alternative is independent-certificate stability: prove directly that every accepted certificate at the relevant query has the same observation, or observations within a prescribed diameter.  Then least selection may be performed independently.  This is developed in \cref{sec:observational-descent}.

\section{Stable observers and finite observational descent}
\label{sec:observational-descent}

The internal witness chosen by a finite search is often not mathematically canonical.  Roots can permute, eigenspaces can have many bases, charts can overlap, and local representatives can differ.  FINER does not force canonical witnesses.  It proves that the requested \emph{observation} is independent of the witness, invariant modulo an explicit equivalence, or stable within an explicit error.

\subsection{Three observer interfaces}

Let \(C(q,c)=1\) mean that \(c\) is an accepted certificate, and let \(O(q,c)\) be an observable code.

\begin{definition}[Exact observer]
An exact observer satisfies
\begin{equation}
 C(q,c)=1\wedge C(q,c')=1
 \;\to\;
 O(q,c)=O(q,c').
 \label{eq:exact-observer}
\end{equation}
\end{definition}

\begin{definition}[Quotient observer]
Given a reflected equivalence relation \(E\), a quotient observer satisfies
\begin{equation}
 C(q,c)=1\wedge C(q,c')=1
 \;\to\;
 E(O(q,c),O(q,c'))=1.
 \label{eq:quotient-observer}
\end{equation}
No quotient type need be constructed; downstream predicates are required to be congruent under \(E\).
\end{definition}

\begin{definition}[Metric observer]
Given a rational-valued primitive-recursive pseudo-distance \(d\), a metric observer has a diameter bound \(\varepsilon(q)\) and satisfies
\begin{equation}
 C(q,c)=1\wedge C(q,c')=1
 \;\to\;
 d(O(q,c),O(q,c'))\leq\varepsilon(q).
 \label{eq:metric-observer}
\end{equation}
\end{definition}

\begin{theorem}[Finite choice erasure]
\label{thm:choice-erasure}
Assume an accepted certificate exists.  Let \(c^*=\Sel(q)\).  Then:
\begin{enumerate}[label=(\roman*),leftmargin=2.3em]
\item under \eqref{eq:exact-observer}, every accepted certificate has observation \(O(q,c^*)\);
\item under \eqref{eq:quotient-observer}, every accepted observation is \(E\)-equivalent to \(O(q,c^*)\);
\item under \eqref{eq:metric-observer}, every accepted observation lies in the \(\varepsilon(q)\)-ball around \(O(q,c^*)\).
\end{enumerate}
\end{theorem}

\begin{proof}
By \cref{thm:selected-certificate}, \(c^*\) is accepted.  Instantiate the corresponding pairwise observer property with \((c,c^*)\).  No choice or coherence argument is needed.
\end{proof}

\subsection{Independent-certificate regularity}

Suppose each precision \(k\) has a nonempty accepted set \(C_k\), and \(O_k(c)\) is a rational approximation.  A particularly robust name interface is
\begin{equation}
 C_k(c)=1\wedge C_\ell(c')=1\wedge k\leq\ell
 \;\to\;
 d(O_k(c),O_\ell(c'))\leq 2^{-k}.
 \label{eq:independent-stability}
\end{equation}
Define \(a_k=O_k(\Sel(k))\).  The selected certificates need not refine one another: \eqref{eq:independent-stability} immediately gives
\[
 d(a_k,a_\ell)\leq2^{-k}\qquad(k\leq\ell).
\]
Thus \(a\) is a regular Cauchy name by a direct object-language theorem.  This is stronger than adjacent compatibility and should be used whenever available.

If only adjacent relational compatibility is known, a bounded consumer should use backward viability to construct the finite compatible segment it actually needs.  An infinite coherent branch must not be inferred without an additional producer.

\subsection{Finite observational descent}

Local-to-global arguments often construct local witnesses on a finite cover and then invoke quotient, sheaf, or gluing language.  When the final theorem only asks for a stable observation, the global witness can be avoided.

Let \(G=(V,E)\) be a finite connected graph.  Each vertex \(v\) has a local observation \(o_v\).  Supply a rooted spanning tree \(T\subseteq E\).  For each tree edge \(e=(u,v)\), suppose one has either exact agreement or a metric estimate.

\begin{theorem}[Exact observational descent]
\label{thm:exact-descent}
If \(o_u=o_v\) on every edge of a rooted spanning tree, then \(o_v=o_r\) for every vertex \(v\), where \(r\) is the root.
\end{theorem}

\begin{proof}
Each non-root vertex has a parent of smaller depth.  Induct on the depth.  At depth zero the statement is reflexivity.  If \(v\) has parent \(p(v)\), the edge hypothesis gives \(o_v=o_{p(v)}\), and the induction hypothesis gives \(o_{p(v)}=o_r\).  Equality transitivity yields \(o_v=o_r\).
\end{proof}

The object proof packs parent and depth arrays, checks that every non-root parent edge is present and decreases depth, and applies quantifier-free induction to a fold over depth or to the maximum depth.

\begin{theorem}[Metric observational descent]
\label{thm:metric-descent}
Let \(d\) satisfy reflexivity and the triangle inequality on the represented observations.  Suppose each rooted tree edge \(e\) has
\[
 d(o_u,o_v)\leq\varepsilon_e.
\]
Then for each vertex \(v\),
\begin{equation}
 d(o_v,o_r)\leq
 \sum_{e\in\operatorname{path}_T(v,r)}\varepsilon_e.
 \label{eq:tree-error}
\end{equation}
In particular, every local observation lies in the ball of radius
\[
 E_T=\max_{v\in V}\sum_{e\in\operatorname{path}_T(v,r)}\varepsilon_e
\]
around the root observation.
\end{theorem}

\begin{proof}
Induct on depth.  If \(v=r\), use reflexivity and the empty sum.  Otherwise let \(p=p(v)\).  The triangle inequality gives
\[
 d(o_v,o_r)\leq d(o_v,o_p)+d(o_p,o_r).
\]
Apply the edge estimate to the first term and the induction hypothesis to the second.  Their sum is exactly the path sum for \(v\).  The maximum conclusion follows by finite comparison.
\end{proof}

\subsection{Quotient descent}

The exact theorem extends verbatim to a reflected equivalence relation \(E\) with object proofs of reflexivity and transitivity.  Edge hypotheses \(E(o_u,o_v)=1\) imply \(E(o_v,o_r)=1\).  A downstream predicate \(P\) may be transported from \(o_r\) to every \(o_v\) whenever it has a checked congruence theorem
\[
 E(x,y)=1\wedge P(x)\to P(y).
\]
This realizes quotient reasoning without constructing quotient objects inside the arithmetic language.

\subsection{Why a spanning tree is enough}

A full finite cover may have many redundant overlaps.  The spanning tree reduces the proof to \(|V|-1\) compatibility edges.  Non-tree edges can remain useful as independent validation but are unnecessary for descent.  A tree itself can be supplied as a certificate and checked by:
\begin{itemize}[leftmargin=2em]
\item one distinguished root;
\item a parent for every other vertex;
\item parent edges belong to the cover graph;
\item depth decreases along parents;
\item all depths are below \(|V|\).
\end{itemize}
These conditions imply connectedness and acyclicity by finite descent.  They are primitive-recursive and quantifier-free after reflection.

\subsection{Applications of observational descent}

The theorem directly addresses several recurring constructions.

\begin{description}[style=nextline,leftmargin=2.4em]
\item[Finite sheaf gluing.]  If the requested output is an invariant of local sections and overlap compatibility makes that invariant equal, use exact descent; do not build a global section unless the theorem needs one.

\item[Charts and coordinate choices.]  Local coordinate formulas may differ, but a scalar, projective, or tensorial observation can descend through checked transition identities.

\item[Root and eigenvalue clusters.]  Local orderings may permute.  Symmetric polynomials, cluster multisets, characteristic-polynomial factors, and spectral projectors are quotient observations.

\item[Approximate geometry.]  Local reconstructions with overlap error bounds yield a root reconstruction and an explicit global error budget by \cref{thm:metric-descent}.
\end{description}

The method does not prove the existence of local data or overlap bounds; those remain producer obligations.  It proves that once finite local certificates exist, an unnecessary infinitary gluing step can often be eliminated.

\section{Representing mathematical concepts by finite observations}
\label{sec:representations}

FINER is not tied to one mathematical area.  Its generality comes from a strict rule: an ``object'' is used only through a primitive-recursive family of finite observations and explicit laws relating those observations.  This section gives reusable interfaces.

\subsection{Natural and integer algebra}

Natural-number algebra is native.  Signed integers are represented by differences \((a,b)\), denoting \(a-b\), with equality
\[
  (a,b)\sim(c,d)\quad\Longleftrightarrow\quad a+d=b+c.
\]
The existing \texttt{Core.Integers} development implements this representation and reduces polynomial identities to ring normalization in the metatheory followed by checked object derivations.  Rational numbers can be represented by signed numerators and positive denominators; equality and order are cleared by cross multiplication.  Division is never used without a positive denominator or a quantitative apartness guard.

\subsection{Finite algebraic structures}

A finite group, ring, category, graph, simplicial complex, or incidence structure is a tuple of bounds and operation/relation tables.  Its axioms are bounded formulas and therefore have reflected bits.  Theorems about one supplied finite structure reduce to finite table checks plus bounded induction.

A family with size parameter \(n\) uses primitive-recursive table accessors and a variable-length fold.  A general theorem such as associativity of a generated operation still requires an object proof; merely checking finitely many instances does not establish a universal family theorem.  Conversely, a theorem about a particular finite presentation should not be burdened with an abstract infinite structure.

\subsection{Regular rational names}

A real-like quantity is represented by a primitive-recursive rational approximation \(a(k)\) and a regularity theorem such as
\begin{equation}
  |a(k)-a(\ell)|\leq 2^{-k}
  \qquad(k\leq\ell).
  \label{eq:regular-name}
\end{equation}
All powers of two and cleared rational inequalities are primitive recursive.  The name is not a completed set-theoretic real inside \iPRA; it is a family of finite claims.  Equality at precision \(k\) means a specified rational error bound, not literal code equality.

Operations on names require precision transformers.  For addition, precision \(k+1\) for both inputs suffices for precision \(k\) of the sum.  Multiplication additionally needs magnitude bounds.  Reciprocal needs a positive apartness margin.

\begin{principle}[Guarded totalization]
Every partial operation is implemented as a total primitive-recursive program with a default output, together with a guard and a correctness theorem under that guard.  The guard must contain the numerical information needed by the error proof.
\end{principle}

For reciprocal, \(x\neq0\) is generally insufficient as finite numerical data.  A usable guard is
\begin{equation}
  |x|\geq 2^{-m}
  \label{eq:apartness-margin}
\end{equation}
for an explicit \(m\).  From approximations to \(x\) finer than \(m+k+O(1)\), rational interval arithmetic produces a \(2^{-k}\)-approximation to \(1/x\).  The demand transformer records this dependence.

\subsection{Functions and operators}

A represented function is a fixed primitive-recursive evaluator
\[
  F(k,x_{\leq d(k)})
\]
together with an error or modulus theorem.  There is no higher-type oracle in the object language; the theorem-specific evaluator reads a finite code or a finite table.  Uniform continuity on a represented compact domain is expressed by a primitive-recursive modulus \(\omega(k)\), and a finite-net theorem uses that modulus explicitly.

When an application genuinely needs an input function as a parameter, use one of three safe interfaces:
\begin{enumerate}[label=(\roman*),leftmargin=2.3em]
\item a fixed function symbol selected at the metalevel;
\item a bounded table sufficient for the current query;
\item a code for a restricted program class with a separately verified clocked evaluator.
\end{enumerate}
An unrestricted total universal primitive-recursive evaluator is unavailable by the diagonal barrier in \cref{sec:design}.

\subsection{Finite measures and probability laws}

A finite rational probability law is a list of natural weights \(w_i\) with positive total \(W\), and an observable is a rational table \(X_i\).  Expectations are represented without division by the pair
\[
  S=\sum_i w_iX_i,\qquad W=\sum_i w_i,
\]
or by cleared identities.  Linearity, monotonicity, indicators, variance identities, Markov and Chebyshev inequalities, finite conditioning, and finite independence are bounded-sum theorems.

A countable or continuous law is not imported wholesale.  It is supplied through finite quantizations \(\mu_k\), tail bounds, moment bounds, anti-concentration data, or an integration modulus.  Every limit claim is reduced to a finite query whose error budget names the required residual data.  This prevents a semantic measure-theory theorem in Lean from being mistaken for an \iPRA{} derivation.

\subsection{Compact spaces and finite nets}

A compact metric presentation consists, for each mesh parameter \(k\), of:
\begin{itemize}[leftmargin=2em]
\item a finite list of points \(N_k\);
\item a coverage certificate: every represented point lies within \(\delta_k\) of some element of \(N_k\);
\item primitive-recursive distance approximations with error bounds;
\item refinement or comparison data when cross-scale arguments need them.
\end{itemize}
Abstract compactness is replaced by the finite coverage formula actually consumed by the proof.  A minimum or finite subcover is then found by bounded search.  The production of \(N_k\) and its coverage proof may be trivial for boxes and finite complexes, or a serious residual for a general represented space.

\subsection{Topology by finite complexes and carriers}

A topological observation at precision \(k\) is a finite simplicial or cubical complex, together with carrier maps linking consecutive scales.  Chains, boundaries, cycles, homologies over a fixed finite coefficient ring, and simplicial maps are finite arrays.  A map that is not uniquely discretized may return a finite cloud; a strong carrier condition makes every selection simplicial, and contiguity yields an explicit finite chain homotopy.

This interface supports finite homology observations and quantitative separation arguments.  It does not by itself prove that a given continuous space has a suitable tower or that the tower completely captures an infinitary topological invariant.  Those are producer and comparison theorems.

\subsection{Algebraic and Diophantine geometry}

Polynomial systems are exact finite data.  Localizations, ideal-membership witnesses, Macaulay matrices, resultants, and bounded-height point searches are finite algebraic certificates.  Quotient equality should be represented by explicit witnesses, such as an ideal combination or a killing exponent, rather than by a semantic quotient type.

For rational or integral points, a verifier can check a point and can enumerate all points below a supplied height.  Completeness requires a height bound or a complete-list theorem.  Cardinality finiteness alone is not an enumeration bound; this issue is treated in \cref{sec:difficulties}.

\subsection{Observational setoids}

Code equality is usually too strict.  A represented concept should carry a reflected equivalence relation \(E\) and congruence proofs for every operation and observer used by the theorem.  This is an \emph{observational setoid}.  One works with codes throughout, transports predicates along \(E\), and uses quotient observers as in \cref{sec:observational-descent}.  There is rarely a need to construct an actual quotient sort inside the single-sorted arithmetic language.

\subsection{A uniform interface record}

The following metalevel data cover most applications:
\begin{align*}
 &\mathsf{Query},\quad \mathsf{CodeBound}(q),\quad \mathsf{Valid}(q,c),\\
 &\mathsf{Observe}(q,c),\quad \mathsf{Equiv}(q,o,o'),\\
 &\mathsf{Refine}(q,q',c,c'),\quad \mathsf{Error}(q),\\
 &\mathsf{ProducerResiduals}(q),\quad \mathsf{ConsumerGoal}(q,o).
\end{align*}
After specialization, every item is a fixed primitive-recursive symbol or formula, and every law is either an inherited/derived object sequent or an explicit residual.  This record is not meant to be one monolithic Lean type imported by every theorem; it is a checklist for designing theorem-specific interfaces.

\section{Explicit residual producers}
\label{sec:residuals}

A verifier theorem and a producer theorem solve different problems.  The verifier proves that finite evidence has a stated consequence.  The producer proves that suitable evidence exists.  In elementary arithmetic both may be easy; in effective analysis or Diophantine geometry the producer may contain almost all of the depth.  FINER keeps the distinction visible in the final sequent.

\subsection{Residual normal form}

A compiled theorem has the form
\begin{equation}
  B_1(q),\ldots,B_r(q),\Gamma(q)
  \;\vdash[\PRAAx]\;\Phi(q),
  \label{eq:residual-normal-form}
\end{equation}
where each \(B_i\) is a concrete object formula, not a prose promise.  Typical residuals are:
\begin{align*}
 &\exists c<M(q)\;[\CertCheck(q,c)=1], &&\text{accepted certificate existence},\\
 &\forall x\,[P_q(x)\to h(x)\leq H(q)], &&\text{height cutoff},\\
 &0<\delta(q)\leq |a-b|, &&\text{separation gap},\\
 &\mathsf{Covers}(N_k,X,\varepsilon_k)=1, &&\text{finite-net coverage},\\
 &\forall n\geq N(k)\;E(n,k)\leq 2^{-k}, &&\text{tail modulus},\\
 &\mathsf{RepresentativeComplete}(L,q)=1, &&\text{covering representatives},\\
 &\mathsf{LocalEstimate}(v,d_v)=1. &&\text{node-specific quantitative estimate}
\end{align*}
The formulas may be universally closed or left open in parameters according to how they are consumed.  They should use the exact encoding, positivity convention, and error metric of the target theorem.

\subsection{Closing residuals by checked simultaneous cut}

\begin{theorem}[Residual closure]
\label{thm:residual-closure}
Suppose
\[
  h:B_1,\ldots,B_r,\Gamma\vdash\Phi
\]
and, under a context \(\Delta\), one has derivations of every formula in \(B_1,\ldots,B_r,\Gamma\).  Then \(\Delta\vdash\Phi\).
\end{theorem}

\begin{proof}
This is exactly the existing theorem \texttt{Prf.trans}.  It simultaneously replaces every hypothesis of \(h\) by its derivation under \(\Delta\).
\end{proof}

In Lean, the closing step has the shape
\begin{lstlisting}[language={}] 
exact Prf.trans compiled fun psi hpsi => producerProof psi hpsi
\end{lstlisting}
No bespoke cut theorem or compiler-specific logic is needed.

\subsection{Residual dependency graphs}

Residuals themselves may consume other quantitative data.  For example, a reciprocal estimate consumes an apartness margin; a root-separation producer consumes coefficient heights; a finite-cover producer consumes a compactness modulus.  Their finite dependencies form the demand graph of \cref{sec:demand}.  The least demand closure fixes all parameters before residual formulas are generated.

A residual graph is permitted to contain a genuine mathematical cycle only if it has a finite monotone capped interpretation.  Otherwise the interface is ill-founded and must be redesigned.  The overflow bit prevents arbitrary caps from being mistaken for a solution.

\subsection{Producer classes}

It is helpful to classify producers by how they can be discharged.

\begin{description}[style=nextline,leftmargin=2.8em]
\item[Arithmetic producer.]  The residual follows from existing checked arithmetic, perhaps after a new short derivation.  Example: a bound on a finite sum.

\item[Algorithmic producer.]  A primitive-recursive construction yields a certificate and a trace proof.  Example: a rational bisection interval for a monotone polynomial.

\item[Input contract.]  The theorem is intentionally conditional on supplied quantitative data.  Example: a continuous-function code together with a modulus.

\item[Deep mathematical producer.]  Closing the residual is the substantive theorem.  Example: an effective height bound for all rational points on a general curve.

\item[Obstructed producer.]  The requested object cannot be uniformly selected with the stated representation.  The consumer must be reformulated to use a relation, quotient, cluster, or guard.  Example: a globally exact ordering of colliding roots.
\end{description}

A good report states which class each residual belongs to.  Calling a deep or obstructed producer ``future engineering'' is misleading.

\subsection{Producer/verifier equivalences}

In several cases one can prove an exact equivalence that clarifies the missing ingredient.

\begin{theorem}[Height bound and complete bounded list]
\label{thm:height-list}
Let \(P(q,x)\) be a decidable predicate and \(h(x)\) a primitive-recursive height.  The following data are interconvertible by primitive-recursive constructions with object proofs:
\begin{enumerate}[label=(\roman*),leftmargin=2.3em]
\item a bound \(H(q)\) such that \(P(q,x)\to h(x)\leq H(q)\);
\item a finite list containing exactly all \(x\) satisfying \(P(q,x)\), provided each code with height at most \(H(q)\) lies below a primitive-recursive enumeration bound \(E(H(q))\).
\end{enumerate}
\end{theorem}

\begin{proof}
From a height bound, enumerate all codes \(x<E(H(q))\), retain those with \(P(q,x)\), and remove duplicates by least occurrence.  Decidability and bounded search prove soundness and completeness.  Conversely, from a complete list \(x_0,\ldots,x_{n-1}\), take the maximum of the finitely many heights.  Completeness shows every solution's height is below that maximum.  All list operations and maxima are primitive recursive and their correctness follows by quantifier-free induction.
\end{proof}

The enumeration-bound premise is often implicit in a coding convention; it should be stated.  A cardinality bound alone is insufficient: the singleton predicate \(x=N\) has cardinality one for every arbitrarily large \(N\).

\subsection{Residual minimization}

The compiler should simplify residual lists without strengthening them.  Safe simplifications include:
\begin{itemize}[leftmargin=2em]
\item syntactic duplicate removal;
\item replacing a residual by a checked equivalent normal form;
\item discharging it immediately from inherited arithmetic;
\item fusing several finite witnesses into one master certificate;
\item replacing a family of local bounds by its checked finite maximum.
\end{itemize}
Unsafe simplifications include replacing an explicit modulus by qualitative convergence, replacing a positive margin by mere inequality, or replacing a height cutoff by a cardinality assertion.  The theorem ledger in the archive records residuals at their exact strength.

\section{The FINER compiler theorem}
\label{sec:compiler}

This section assembles the preceding components into a single metatheorem.  The theorem is deliberately stated for a \emph{fixed source graph}.  Its Lean formalization is a verified compiler that constructs object syntax and derivations, not an object-level interpreter.

\subsection{Specialized source graph}

A source graph \(G\) contains:
\begin{enumerate}[label=(\roman*),leftmargin=2.3em]
\item finitely many nodes of the six classes in \cref{sec:design};
\item primitive-recursive query encoders and node-local programs;
\item monotone capped demand transformers;
\item local preconditions, postconditions, and residual formulas;
\item local object derivations proving node soundness at arbitrary terms;
\item a terminal formula \(\Phi_G(q,o)\).
\end{enumerate}
Every syntax tree in this list is ordinary Lean data whose leaves are fixed \texttt{PR} symbols or already constructed object derivations.

\subsection{Compilation passes}

The compiler performs the following finite metalevel passes.

\begin{enumerate}[label=\textbf{P\arabic*.},leftmargin=3.2em]
\item \textbf{Demand extraction.}  Collect every node's computation and justification demands into a monotone capped transformer \(F_G\).

\item \textbf{Demand closure.}  Generate the primitive-recursive iteration of \(d\mapsto d\vee F_G(d)\), its rank bound, and the least-closure derivations of \cref{thm:demand-stabilization,thm:demand-least}.  Add overflow zero as a residual or a checked consequence.

\item \textbf{Specialization.}  Substitute the closed demand vector into every node.  Clocked loops and relational interactions now have explicit finite horizons and bounds.

\item \textbf{Certificate layout.}  Allocate mixed-radix fields for all non-recomputed evidence and define the master bound \(M_G(q)\), decoder, checker \(C_G(q,c)\), and observer \(O_G(q,c)\).

\item \textbf{Checker reflection.}  Compile all bounded checks to a normalized bit and construct the reflection derivations.

\item \textbf{Local proof schedule.}  Unroll finite cyclic execution into certified iterations or viability tables, then order the resulting local sequents.  Compose them by substitution and simultaneous cut to obtain
\begin{equation}
  C_G(q,c)=1,\;\vec B_G(q),\;\Gamma_G(q)
  \vdash[\PRAAx]
  \Phi_G(q,O_G(q,c)).
  \label{eq:generic-certificate-soundness}
\end{equation}

\item \textbf{Selection.}  Use checked bounded search to define \(c_G^*(q)\).  From the certificate-existence residual, derive \(C_G(q,c_G^*)=1\) and substitute it into \eqref{eq:generic-certificate-soundness}.

\item \textbf{Residual closure.}  Leave \(\vec B_G\) in the public endpoint or close selected members with supplied producer derivations using \texttt{Prf.trans}.
\end{enumerate}

\begin{theorem}[FINER compilation]
\label{thm:finer-compilation}
For every well-formed fixed source graph \(G\) whose local node derivations are values of the ordinary \(\Prf\;\PRAAx\) calculus, the passes above effectively construct primitive-recursive terms \(M_G,C_G,O_G,c_G^*\), an explicit finite residual list \(\vec B_G\), and an object derivation
\begin{equation}
  \vec B_G(q),\Gamma_G(q)
  \;\vdash[\PRAAx]\;
  \Phi_G(q,O_G(q,c_G^*(q))).
  \label{eq:finer-compiled-theorem}
\end{equation}
No principle stronger than the axioms \(\PRAAx\) and intuitionistic first-order logic is introduced.
\end{theorem}

\begin{proof}
We verify each pass.

\paragraph{Passes P1--P2.}
The global demand transformer is assembled from fixed primitive-recursive node transformers using componentwise maximum and clipping.  Monotonicity is proved from the node monotonicity derivations.  The finite demand-closure theorems provide a least prefixed point and a quantifier-free proof of adequacy.

\paragraph{Pass P3.}
Object-language substitution is admissible for every derivation.  Substituting the demand terms into node-local proofs therefore preserves their conclusions.  Metalevel specialization of fixed syntax creates fixed \texttt{PR} symbols; no universal evaluator is used.

\paragraph{Pass P4.}
Finite products and sums define the mixed-radix bound and decoders.  The correctness lemmas of \cref{sec:finite-data} prove that every field read by the checker is the stored field and that every code below the bound has a canonical field vector.

\paragraph{Pass P5.}
The checker is a fixed bounded formula.  Apply \cref{thm:boolean-reflection}.  Master acceptance then yields every local check formula by propositional elimination and bounded witness extraction.

\paragraph{Pass P6.}
Exact nodes rewrite by their defining equations.  Boolean and search nodes use reflection and the checked search theorems.  Iteration nodes use \cref{thm:trace-reflection}.  Relational nodes use \cref{thm:least-path}; observation nodes use \cref{thm:choice-erasure} or the descent theorems.  Every local result is an ordinary sequent.  The finite schedule is composed with substitution, weakening, implication elimination, and simultaneous cut, all already admissible in \(\Prf\).

\paragraph{Pass P7.}
The residual list includes an accepted-certificate existence formula, unless a producer derivation has already discharged it.  Apply \cref{thm:selected-certificate} to obtain master acceptance at \(c_G^*\).  Instantiate \eqref{eq:generic-certificate-soundness} with this term and cut in acceptance.

\paragraph{Pass P8.}
Retaining residuals gives \eqref{eq:finer-compiled-theorem}.  Supplying their derivations and applying \cref{thm:residual-closure} gives an unconditional theorem under the remaining context.  Every logical operation used is a constructor or proved structural theorem of the existing ordinary calculus.
\end{proof}

\subsection{Soundness of the metalevel compiler}

A Lean implementation can make \cref{thm:finer-compilation} literal: define an inductive type of well-formed node packages whose fields include their local \texttt{Prf} derivations.  The compiler recursively returns a package containing terms, formulas, residuals, and a final \texttt{Prf}.  Because \texttt{Prf} is an inductive proposition checked by Lean's kernel, a bug in the metalevel compiler cannot create a proof unless it also constructs a valid derivation value.

The compiler's semantic reference algorithms may still be tested independently, but those tests are not trusted by the object theorem.  Generated terms and derivations are the certificate.

\subsection{Open contexts and portability}

The theorem is intentionally stated under \(\Gamma_G(q)\), not only for closed sentences.  Local hypotheses may encode an input contract, algebraic assumptions, or a theorem-specific antecedent.  This makes the result reusable at arbitrary terms.

The current public Foundation bridge primarily exposes closed-sentence transport.  A compiled open theorem can be universally closed before transport, or it can wait for the planned open-sequent bridge.  The canonical maintained proof should remain the ordinary \(\Prf\;\PRAAx\) endpoint; backend-specific statements are corollaries.

\subsection{What is not claimed}

The compiler theorem does not assert that every informal mathematical proof can be specialized automatically.  It assumes:
\begin{itemize}[leftmargin=2em]
\item a finite-query formulation;
\item fixed primitive-recursive node programs;
\item local object derivations;
\item monotone capped demand data;
\item explicit residual producer formulas.
\end{itemize}
Finding these ingredients may require substantial mathematics.  FINER's contribution is to make their composition uniform and the remaining gap exact.

\section{Resolution matrix for recurring difficulties}
\label{sec:difficulties}

This section revisits the obstacles repeatedly identified across the proof packets.  ``Resolved'' means that FINER gives a generic finite proof schema.  ``Conditional'' means that the proof is generic once named quantitative data are supplied.  ``Reformulated'' means that the stronger requested interface is not constructively stable and should be replaced.  ``Open producer'' means that the finite verifier is understood but the mathematical existence theorem remains deep.

\subsection{Ambient evidence versus an object proof}

\paragraph{Difficulty.}
A Mathlib theorem, executable Python function, or finite test is cited as if it were an \iPRA{} proof.

\paragraph{Resolution.}
The public endpoint must have literal type \(\Gamma\vdash[\PRAAx]\varphi\).  Host theorems may generate terms, normalize polynomial identities, or test examples, but every imported fact is represented by a local object sequent.  FINER's node package includes the sequent as data.  Status: \textbf{resolved architecture}; theorem-specific adapters remain work.

\subsection{Bounded quantifiers inside induction}

\paragraph{Difficulty.}
An invariant contains \(\forall i<N\) or \(\exists i<N\), so it is not eligible for quantifier-free induction.

\paragraph{Resolution.}
Compile the bounded formula to a normalized bit and induct on the equation that the bit is one.  Finite sums, products, and least search supply reflection and witness extraction.  Status: \textbf{resolved}; the missing deliverable is a checked generic reflection library.

\subsection{All precisions maintained inside one loop}

\paragraph{Difficulty.}
A construction attempts to carry an entire Cauchy name or every future precision through a finite induction.

\paragraph{Resolution.}
Fix the terminal query first.  Run reverse demand propagation, solve finite cycles by \cref{thm:demand-stabilization}, and execute only the finite specialization required for that query.  Quantify over the query outside the trace theorem.  Status: \textbf{resolved schema}.

\subsection{Cyclic precision dependencies}

\paragraph{Difficulty.}
Node \(A\) needs precision from \(B\), while \(B\) needs precision from \(A\); an acyclic reverse pass fails.

\paragraph{Resolution.}
Use a monotone capped transformer on a finite product of chains.  The iteration \(d\mapsto d\vee F(d)\) reaches the least prefixed point after at most the sum of capacities.  Check an overflow bit to ensure caps are adequate.  Status: \textbf{resolved for finite monotone capped cycles}.  Nonmonotone or unbounded cycles require redesign or an explicit clocked search.

\subsection{Hidden finite choice}

\paragraph{Difficulty.}
Separate existential witnesses at multiple stages are assumed to fit together.

\paragraph{Resolution.}
Pack all evidence into one master certificate, or compute a backward viability table and select the least successful path.  Pointwise existence is never used as compatibility.  Status: \textbf{resolved for bounded interactions}.

\subsection{Classical case analysis}

\paragraph{Difficulty.}
A finite proof invokes excluded middle or contradiction without recovering constructive evidence.

\paragraph{Resolution.}
For fixed bounded formulas, Boolean reflection proves decidability and supplies success/failure certificates.  Classical propositional reasoning over those atoms is simulated by finite case splits.  For unbounded formulas, a negative translation is sound, but recovery of \(\Phi\) from \(\neg\neg\Phi\) requires a stability theorem; positive existential witnesses require explicit certificates.  Status: \textbf{resolved for bounded decidable reasoning; conditional for stable unbounded conclusions}.

\subsection{Exact sign and branch selection}

\paragraph{Difficulty.}
An algorithm is asked to choose an exact sign at zero, one ordered root through a collision, or one branch of a multivalued analytic function without separation data.

\paragraph{Barrier.}
A finite-query algorithm examines only finitely much of a Cauchy input.  At the zero name, choose a stage beyond every query it makes.  There are positive and negative regular names agreeing with zero on all queried approximations and separating only later.  The same output cannot be a correct exact sign for both.  Thus no such finite continuous selector is justified without a gap.

\paragraph{Resolution.}
Return overlapping branches, a finite relation, a root multiset or cluster, a projective direction, or another quotient observation.  Recover a unique label only under a positive separation residual.  Status: \textbf{reformulated}; gap-guarded selection is conditional and finite.

\subsection{Eigenvectors and rotation axes at degeneracy}

\paragraph{Difficulty.}
A repeated eigenvalue is assigned a canonical eigenvector, or the identity rotation is assigned a canonical axis.

\paragraph{Resolution.}
Use invariant subspaces, spectral projectors, Gram matrices, quaternions modulo sign, or projective relations.  A vector or axis may be selected when a spectral gap, normalization coordinate, or distance from the degenerate locus is supplied.  Observational descent handles chart changes.  Status: \textbf{reformulated}; exact vector selection is gap-conditional.

\subsection{Partial operations from mere nonzeroness}

\paragraph{Difficulty.}
A quantitative reciprocal, normalization, pivot, or division algorithm assumes only \(x\neq0\).

\paragraph{Resolution.}
Totalize the operation and require a positive apartness margin, lower denominator bound, pivot bound, or determinant bound.  Reverse demand propagation computes the necessary input precision.  Status: \textbf{conditional resolution}.  Mere negated equality is logically weaker than the numerical resource consumed by an error estimate.

\subsection{Abstract compactness}

\paragraph{Difficulty.}
A finite subcover, minimizer, or mesh representative is selected from an abstract compactness theorem.

\paragraph{Resolution.}
Use a primitive-recursive finite net, coverage checker, and modulus.  Bounded search selects the finite witness.  If the represented space does not already provide a net and coverage theorem, they are residual producers.  Status: \textbf{conditional resolution}.

\subsection{Metastability promoted to a convergence rate}

\paragraph{Difficulty.}
A finite-window stability witness is treated as a global Cauchy modulus.

\paragraph{Barrier.}
Any computation that inspects only a finite prefix can be fooled by a sequence that agrees with a constant sequence on that prefix and changes afterward.  More concretely, after a proposed rate \(R(k)\) has been computed from finite window data, define a binary sequence that remains zero through every inspected index and becomes one after \(R(k)\).  The same finite evidence cannot certify the requested tail bound for both sequences.

\paragraph{Resolution.}
Retain the metastable statement for finite consumers.  To construct a regular name or global rate, add a tail modulus or a summable adjacent-motion bound as a residual.  Status: \textbf{stronger conversion unavailable in general}.

\subsection{Pointwise approximants mistaken for a coherent name}

\paragraph{Difficulty.}
For every precision there exists an approximation, and these are treated as one Cauchy name.

\paragraph{Resolution.}
Use independent-certificate stability \eqref{eq:independent-stability}, or a finite compatible path from backward viability plus a tail producer.  Diameter bounds within each precision alone are not enough; cross-precision motion must be controlled.  Status: \textbf{conditional resolution}.

\subsection{Puiseux and algebraic branches}

\paragraph{Difficulty.}
A uniformly exact Puiseux branch is selected from arbitrary Cauchy coefficients, including collisions and coefficient perturbations.

\paragraph{Resolution.}
For exact algebraic or splitting input, finite factor and Newton-polygon certificates can name branches.  For approximate coefficients, return clusters, truncated expansions with relational supports, or symmetric observations.  Exact labels require separation and exact-presentation residuals.  Status: \textbf{unrestricted selector obstructed; verifier architecture resolved}.

\subsection{Diophantine finiteness versus enumeration}

\paragraph{Difficulty.}
A proof that a solution set is finite, or has cardinality at most \(N\), is treated as an effective cutoff for bounded search.

\paragraph{Resolution.}
Require a height bound or a complete-list theorem.  \cref{thm:height-list} converts between them under an explicit code-enumeration bound.  Cardinality alone gives no location information: singleton solution sets can occur at arbitrary heights.  Status: \textbf{verifier resolved; effective height producer remains deep}.

\subsection{Mordell--Weil and descent representatives}

\paragraph{Difficulty.}
Finite Selmer or covering classes are treated as if rational representatives were automatically available.

\paragraph{Resolution.}
Separate the finite class computation from a representative relation.  A certificate supplies a rational point on a covering, compatibility with the class, local checks, and the descent map.  Backward viability composes only the representatives actually used by a bounded descent.  General generation requires an explicit representative or height producer.  Status: \textbf{verifier architecture resolved; general producer open}.

\subsection{Ineffective analytic constants}

\paragraph{Difficulty.}
A proof uses ``for sufficiently large \(N\)'' or an unspecified constant inside a finite algorithm.

\paragraph{Resolution.}
Move the threshold, constant, and verified estimate into the residual list.  The checker proves the final finite inequality from these values.  A separate mathematical development may later produce them.  Status: \textbf{resolved normal form; producer theorem may be open}.

\subsection{General probability spaces}

\paragraph{Difficulty.}
Measure-theoretic operations hide infinitary approximation and discontinuities.

\paragraph{Resolution.}
Start with finite rational laws.  For larger spaces, supply quantizations, tail moments, integration moduli, and anti-concentration data.  Discontinuous statistics such as exact quantiles require gap or anti-concentration guards; robust intervals are preferable at atoms.  Status: \textbf{conditional resolution}.

\subsection{Completed topological spaces}

\paragraph{Difficulty.}
A proof manipulates a completed compact space without a finite representation of maps, covers, or homotopies.

\paragraph{Resolution.}
Use finite complexes or nets, strong carriers, finite chain maps, and explicit prism homotopies.  Observational towers may approximate an infinite object, but comparison and completeness theorems are explicit producers.  Status: \textbf{representation resolved; theorem-specific approximation producers remain}.

\subsection{Extensionality and code equality}

\paragraph{Difficulty.}
Different codes of the same mathematical object are required to be literally equal.

\paragraph{Resolution.}
Use observational setoids, operation congruence, finite back-and-forth, and quotient observers.  Exact code normalization is optional and theorem-specific.  Status: \textbf{resolved architecture}.

\subsection{Finite gluing turned into an infinitary quotient}

\paragraph{Difficulty.}
A finite local argument constructs an unnecessary global quotient, sheaf object, or chosen representative.

\paragraph{Resolution.}
Use finite observational descent along a checked spanning tree.  Exact edge compatibility gives exact or quotient descent; metric compatibility gives an explicit accumulated error.  Construct a global witness only when the terminal theorem genuinely mentions it.  Status: \textbf{resolved at observation level}.

\subsection{A concise status table}

\begin{longtable}{@{}p{0.24\textwidth}p{0.46\textwidth}p{0.22\textwidth}@{}}
\toprule
Obstacle & FINER treatment & Status \\
\midrule
\endfirsthead
\toprule
Obstacle & FINER treatment & Status \\
\midrule
\endhead
Bounded logic in induction & Boolean reflection and bit invariants & Resolved schema \\
Cyclic demands & Least finite monotone closure & Resolved under caps \\
Nonunique finite witnesses & Backward viability and least replay & Resolved boundedly \\
Exact labels at collisions & Clusters/quotients or gap guards & Reformulated \\
Partial operations & Totalization plus quantitative guard & Conditional \\
Abstract compactness & Finite net and coverage certificate & Conditional \\
Metastability & Retain finite-window contract & No general rate \\
Diophantine finiteness & Height/list producer & Deep residual \\
Classical bounded proof & Dual finite certificates & Resolved boundedly \\
Unbounded existential proof & Explicit witness producer & Residual \\
Local-to-global observation & Spanning-tree descent & Resolved finitely \\
Universal PR evaluator & Metalevel specialization & Impossible interface \\
\bottomrule
\end{longtable}

\section{Worked application patterns}
\label{sec:applications}

The examples below are not substitutes for theorem-specific developments.  They show how broad mathematical constructions factor through the common kernel and where their residual boundaries lie.

\subsection{Exact rational bisection for \(y^3+y=x\)}
\label{sec:cubic-example}

Fix a rational input \(x\) and rationals \(L_0\leq U_0\) satisfying
\begin{equation}
  f_x(L_0)\leq0\leq f_x(U_0),
  \qquad f_x(y)=y^3+y-x.
  \label{eq:cubic-bracket}
\end{equation}
All rational statements are represented by cleared integer inequalities.

\begin{lemma}[Strict monotonicity]
If \(u<v\), then \(f_x(u)<f_x(v)\).
\end{lemma}

\begin{proof}
Factor the difference:
\[
 f_x(v)-f_x(u)
 =(v-u)(v^2+uv+u^2+1).
\]
The first factor is positive.  The second is positive because
\[
 v^2+uv+u^2
 = (v-u)^2+3uv
\]
after one rearrangement is not uniformly nonnegative when \(uv<0\), so that form is not ideal.  Instead use
\[
 v^2+uv+u^2
 = \left(v+\frac{u}{2}\right)^2+\frac{3u^2}{4}\geq0.
\]
To avoid division in the object proof, clear denominators:
\[
 4(v^2+uv+u^2)=(2v+u)^2+3u^2\geq0.
\]
Thus \(v^2+uv+u^2+1>0\).  Multiplication by the positive factor \(v-u\) preserves positivity.  The identity is a polynomial identity suitable for the repository's ring tool; positivity uses the checked order library.
\end{proof}

Define the midpoint \(M_t=(L_t+U_t)/2\) and update
\[
 (L_{t+1},U_{t+1})=
 \begin{cases}
   (M_t,U_t),& f_x(M_t)\leq0,\\
   (L_t,M_t),& f_x(M_t)>0.
 \end{cases}
\]
For exact rational arithmetic, store dyadic-scaled endpoints so division by two is a shift in the denominator.  The branch bit is reflected by order; the update is a total primitive-recursive selector.

\begin{theorem}[Bisection invariant]
For every \(t\),
\begin{align}
 L_t&\leq U_t,\label{eq:bisect-order}\\
 f_x(L_t)&\leq0\leq f_x(U_t),\label{eq:bisect-sign}\\
 U_t-L_t&=2^{-t}(U_0-L_0).\label{eq:bisect-width}
\end{align}
\end{theorem}

\begin{proof}
Use quantifier-free induction on \(t\).  The base is \eqref{eq:cubic-bracket}.  For the step, split on the reflected comparison \(f_x(M_t)\leq0\).  In the first branch, the lower sign condition is the branch hypothesis and the upper condition is inherited; in the second, the lower condition is inherited and the upper condition follows from the negated comparison, using decidability to obtain \(0<f_x(M_t)\) and hence \(0\leq f_x(M_t)\).  The midpoint lies between the endpoints by rational arithmetic.  In either branch the width is halved.  All three clauses are equations or order formulas and can be compiled to a quantifier-free conjunction.
\end{proof}

For query \(k\), choose \(t\) so that \(2^{-t}(U_0-L_0)\leq2^{-k}\).  The observation is the interval \([L_t,U_t]\), or its midpoint with radius half the width.  No completed real root or exact branch selector is required.  Strict monotonicity proves any two valid brackets of width \(2^{-k}\) have compatible observations, yielding a regular name.

The accompanying Python certificate records the initial bracket, number of steps, and final bracket; its verifier recomputes all steps exactly with \texttt{Fraction}.  The object formalization should recompute branch bits and endpoints rather than trust them.

\subsection{Finite-net approximate minimization}

Let \(X\) be represented at query \(k\) by a finite net \(N_k\), and let \(f\) have a modulus guaranteeing
\begin{equation}
 d(x,y)\leq\delta_k\quad\Longrightarrow\quad |f(x)-f(y)|\leq\eta_k.
 \label{eq:min-modulus}
\end{equation}
Assume coverage:
\begin{equation}
 \forall x\in X\;\exists y\in N_k\;d(x,y)\leq\delta_k.
 \label{eq:net-coverage}
\end{equation}
Finite search selects \(y^*\in N_k\) with the least approximated value, including enough evaluation precision to make comparison error at most \(\rho_k\).

\begin{theorem}[Approximate minimum from a certified net]
For every represented \(x\in X\),
\begin{equation}
 f(y^*)\leq f(x)+\eta_k+2\rho_k.
 \label{eq:approx-min}
\end{equation}
\end{theorem}

\begin{proof}
Coverage supplies a net point \(y\) within \(\delta_k\) of \(x\).  By \eqref{eq:min-modulus}, \(f(y)\leq f(x)+\eta_k\).  The selected approximate value is no larger than the approximate value at \(y\).  Replacing both approximations by true values costs at most \(\rho_k\) on each side, giving \eqref{eq:approx-min}.
\end{proof}

The generic verifier and search are resolved.  The residuals are exactly the finite net, coverage theorem, evaluation error, and modulus.  For a rational box they are primitive-recursive producers; for an abstract compact presentation they are input data or deeper theorems.

\subsection{Root clusters and spectral data}

For a monic polynomial with approximate coefficients, exact ordered roots are unstable at collisions.  A FINER interface can instead return:
\begin{itemize}[leftmargin=2em]
\item disjoint rational disks or intervals, each with a certified root count;
\item the multiset of cluster centers and radii;
\item symmetric polynomial observations;
\item a factorization up to a stated coefficient error.
\end{itemize}
The master certificate contains contour, Sturm, resultant, or interval-arithmetic checks appropriate to the coefficient representation.  Internal root orderings are quotient witnesses under permutation.

If a residual proves pairwise cluster separation by \(\delta>0\), least matching between consecutive precisions yields stable labels.  Without that residual, the public theorem remains at the cluster or multiset level.  The same principle applies to matrices: characteristic coefficients and invariant subspaces are unconditional finite observations; individual normalized eigenvectors require a spectral gap and a normalization guard.  Spectral projectors are often the best quotient observer because they remain stable through basis changes within an eigenspace.

\subsection{Diophantine enumeration}

Let \(P(q,x)=0\) be a fixed polynomial or finite Diophantine system and let \(h(x)\) be the chosen height.  Given a residual
\[
  P(q,x)=0\to h(x)\leq H(q),
\]
compute an enumeration bound \(E(H(q))\), scan every code below it, and retain exactly the solutions.  Boolean reflection proves every entry is a solution; the height residual proves every solution occurs.  Duplicate removal uses least occurrence.

The final theorem may state:
\begin{enumerate}[label=(\roman*),leftmargin=2.3em]
\item the output list contains only solutions;
\item every solution appears;
\item a requested aggregate, such as cardinality, maximum height, or a finite sum over solutions, is correct.
\end{enumerate}
The finite enumeration layer is generic.  In Runge-type cases, an explicit estimate may discharge the height residual.  In a general Faltings or Siegel application, the residual identifies the actual missing effective theorem; a finite checker does not close it.

\subsection{Mordell--Weil style finite descent}

A bounded descent consumes:
\begin{itemize}[leftmargin=2em]
\item a finite list of covering classes;
\item for each used class, a finite relation of rational representatives;
\item local solubility and compatibility certificates;
\item a descent map and a strictly decreasing height or rank;
\item a terminal bounded list.
\end{itemize}
Backward viability chooses representatives only along successful branches.  A generic descent theorem proves termination from the decreasing rank and checks reconstruction.  The residual boundary is a representative-completeness theorem and, in general, a height or saturation bound.  This is more accurate than treating a finite cohomological class list as a point list.

\subsection{Finite probability and expectation}

For a finite weighted law \((w_i,X_i)_{i<n}\), define
\[
 W=\sum_{i<n}w_i,
 \qquad
 S=\sum_{i<n}w_iX_i.
\]
The expectation is the rational pair \((S,W)\).  Linearity is a finite-sum identity; monotonicity is a termwise order theorem plus sum monotonicity.  Markov's inequality can be cleared of division:
\[
 a\sum_i w_i[\,X_i\geq a\,]
 \leq \sum_i w_iX_i
\]
after representing indicators by reflected bits.  Chebyshev follows from the same finite Markov argument applied to squared centered values and cleared rational arithmetic.

For an approximated countable law, a query \(k\) supplies a finite truncation and a tail moment bound.  The finite theorem proves the truncated identity and adds the tail error.  The tail bound is a residual or a producer theorem.  Exact quantiles at atoms should be represented by median intervals or guarded by anti-concentration data.

\subsection{Finite topology and carriers}

Suppose a represented map between finite complexes sends each source vertex to a nonempty finite cloud of target vertices.  A strong carrier condition states that the union of clouds over every source simplex lies in one target simplex.  Then every vertexwise selection is simplicial.  Any two selections are contiguous because their combined image of a source simplex lies in the same carrier simplex.  The standard prism formula gives a finite chain homotopy.

Every statement is a bounded array identity.  Backward viability is only needed when clouds at adjacent refinement levels must be compatible.  If the final homology map is invariant under every selection, exact observer erasure eliminates the choices independently.  The residuals are the carrier and refinement certificates, not abstract continuity.

\subsection{Finite local algebra and observational gluing}

Consider a finite affine cover with local module elements \(s_v\) and overlap witnesses.  If the terminal theorem asks only for a polynomial invariant \(I(s_v)\) known to agree on overlaps, choose a spanning tree and apply exact observational descent to obtain one global invariant value.  This avoids constructing a global module element.  If a global element is required, the usual finite equalizer or Cech kernel can itself be a bounded linear-algebra certificate; its existence remains a separate local-to-global producer.

\subsection{Busy Beaver and consistency reductions}

Machine simulations, proof predicates, and bounded searches are especially well suited to FINER:
\begin{itemize}[leftmargin=2em]
\item literal transition tables are finite exact data;
\item run traces are certified iterations;
\item syntax and proof checking are fixed primitive-recursive programs;
\item reduction-specific arithmetic is composed from checked core lemmas;
\item exact machine and theory identifiers remain in the residual/provenance record.
\end{itemize}
The framework does not change the strength of the consistency statement and should not conflate a semantic nonhalting proof, an object-theory equivalence, and a bridge export.  It does provide a disciplined route for sharing arithmetic and trace reflection across such reductions.

\section{Lean formalization plan}
\label{sec:lean}

The formalization should be integrated as a new library beside \texttt{IPRADerivations}, not as copied source inside a packet.  This section gives concrete definitions, theorem shapes, and a dependency order.

\subsection{Imports and namespace}

A minimal first module can import:
\begin{lstlisting}[language={}]
import IPRADerivations.Core.Search
import IPRADerivations.Core.Pairing
import IPRADerivations.Core.MultiSums
\end{lstlisting}
and use
\begin{lstlisting}[language={}]
namespace LeanProofs.IPRA.FINER
open LeanProofs.ConciseIPRA LeanProofs.IPRA Syn
open scoped Fm
\end{lstlisting}
The first objective is to keep this import closure independent of Foundation, Busy Beaver, DiophantineConsistency, and application packets.

\subsection{Kernel wrappers}

The archive includes a small \texttt{KernelWrappers.lean} recording the intended API.  The relevant existing declarations are:
\begin{lstlisting}[language={}]
Prf.subst
Prf.weaken
Prf.trans
Prf.ofNil
LeanProofs.IPRA.inst
LeanProofs.IPRA.ind
Search.search_spec
Search.search_found
Search.search_minimal
\end{lstlisting}
No new structural proof calculus is required.

\subsection{Phase 1: Boolean reflection}

Define a small metalevel syntax whose terms are already \texttt{PTm} values:
\begin{lstlisting}[language={}]
inductive BFormula : (n : Nat) -> Type
| eq  : PTm n -> PTm n -> BFormula n
| lt  : PTm n -> PTm n -> BFormula n
| and : BFormula n -> BFormula n -> BFormula n
| or  : BFormula n -> BFormula n -> BFormula n
| imp : BFormula n -> BFormula n -> BFormula n
| bex : PTm n -> BFormula (n+1) -> BFormula n
| ball : PTm n -> BFormula (n+1) -> BFormula n
\end{lstlisting}
Return a record:
\begin{lstlisting}[language={}]
structure Reflection (n : Nat) where
  bit      : PTm n
  formula  : Fm PTm n
  bit01    : forall {Gamma}, Gamma |- [PRAAx] bit = 0 or bit = 1
  one_iff  : forall {Gamma}, Gamma |- [PRAAx]
               iff (bit = 1) formula
  zero_iff : forall {Gamma}, Gamma |- [PRAAx]
               iff (bit = 0) (neg formula)
\end{lstlisting}
The actual Lean syntax uses the repository notation \(\mathord{\simeq}\), \(\mathbin{\curlywedge}\), and \texttt{Fm.iff}; the pseudo-code above emphasizes the theorem shape.

Implement atoms first.  Then conjunction, disjunction, and implication.  Add bounded existential by a finite sum or a direct adapter to \texttt{Search.mu}; add bounded universal by reflected absence of a counterexample.  At every constructor, prove substitution lemmas for the generated bit and formula.  Acceptance test: a closed example with two nested bounded quantifiers produces a literal \texttt{Prf PRAAx} theorem and passes the repository axiom audit.

\subsection{Phase 2: packed arrays}

Implement mixed-radix arrays or reuse an existing sequence code if its import closure is appropriate.  Required endpoints are:
\begin{align*}
 &\decode_i(\code(\vec d))=d_i,\\
 &c<R\to\code(\decode(c))=c,\\
 &\decode_i(c)<r_i,\\
 &\text{fold equations for sum, maximum, and Boolean all/any.}
\end{align*}
For a first implementation, fix the vector length at the metalevel.  A later variable-length implementation can use \texttt{sumF}, \texttt{prodF}, division, and pairing.

\subsection{Phase 3: trace reflection}

Define a package whose fields are fixed terms and local derivations:
\begin{lstlisting}[language={}]
structure TraceSystem (k : Nat) where
  init : PTm k
  step : PTm (k+2)       -- time, state, parameters
  inv  : Fm PTm (k+2)
  invQF : inv.QF
  base : [] |- [PRAAx] ...
  next : [inv] |- [PRAAx] ...
\end{lstlisting}
Use \texttt{recDef} for the state and \texttt{ind} for \cref{thm:trace-reflection}.  Keep all parameters open.  Prove a separate terminal adapter, so applications can reuse the trace theorem without changing its definition.

\subsection{Phase 4: finite demand closure}

Start with fixed dimension \(m\).  Represent vectors by \texttt{Fin m -> PTm n} at the statement level and by one packed term for iteration.  Define:
\begin{itemize}[leftmargin=2em]
\item componentwise maximum and comparison bits;
\item a fixed-capacity rank term;
\item the inflationary step \(x\vee F(x)\);
\item the trace by \texttt{rec1} or \texttt{rec2}.
\end{itemize}
Prove, in order:
\begin{enumerate}[label=(\roman*),leftmargin=2.3em]
\item coordinate bounds;
\item inflation;
\item strict-step rank increase;
\item persistence after a fixed point;
\item stabilization at \(H=\sum_iC_i\);
\item leastness below an arbitrary prefixed point.
\end{enumerate}
The monotonicity proof of \(F\) should be a field of the system package, not a semantic Lean hypothesis about \(\Nat\).  It must be an object derivation at arbitrary terms.

\subsection{Phase 5: relational viability}

Implement one fixed-horizon version first.  A layer is a fixed bound, an edge is a reflected predicate, and acceptance is reflected.  Generate viability bits backward at the metalevel for a fixed number of layers, then prove path equivalence by structural recursion.  This version already covers many finite proof graphs.

The variable-horizon version packs a viability table and uses a reverse-time iterator.  Its proof should reuse the trace kernel.  The least-path theorem calls \texttt{Search.search\_found} at each forward step.  Add a no-path theorem from viability zero.

\subsection{Phase 6: observers and descent}

Implement exact observer erasure first; it is one instantiation of pairwise invariance.  Then add a reflected equivalence relation with reflexivity/transitivity and congruence.  Metric observers require a rational distance representation, triangle inequality, and finite sums of edge errors.

For observational descent, use a parent/depth certificate rather than a generic graph library.  Prove parent-depth termination and path-error accumulation by quantifier-free induction.  A fixed-size graph version can precede the packed variable-size version.

\subsection{Phase 7: master certificates and compiler packages}

Do not begin with a universal graph language.  Define a \texttt{CompiledQuery} record that stores already generated terms, formulas, local derivations, and residuals.  Prove the selected-certificate theorem and residual closure.  Then add metalevel constructors for exact, reflected, search, trace, relational, and observer nodes.  Each constructor returns a new \texttt{CompiledQuery} and a checked composition proof.

Only after these constructors stabilize should a convenient graph DSL be added.  Its compiler is trusted only through the \texttt{Prf} value it returns.

\subsection{Phase 8: benchmarks}

Four benchmarks exercise different parts of the framework:
\begin{enumerate}[label=(\arabic*),leftmargin=2.3em]
\item rational bisection for \(y^3+y=x\): Boolean branching, trace reflection, regular-name observation;
\item a genuinely cyclic three-coordinate demand system: stabilization and leastness;
\item a finite relational path with a nonviable early choice: backward viability and least replay;
\item a finite cover with rational edge errors: metric observational descent.
\end{enumerate}
The JSON fixtures and Python tests in this archive give exact expected outputs.

\subsection{Phase 9: repository integration}

After the generic library is checked:
\begin{itemize}[leftmargin=2em]
\item migrate one proof packet whose main theorem is already close to a checked derivation;
\item replace its local search or trace code by FINER adapters;
\item record a literal theorem name and axiom audit;
\item preserve the packet report as provenance but make the new derivation the canonical result;
\item expose a Foundation corollary only after the ordinary endpoint is stable.
\end{itemize}
A good first target is a finite algebraic or combinatorial theorem, not a deep analytical producer.  A second target should exercise a residual, such as finite-net minimization under a supplied coverage theorem.

\subsection{Axiom and placeholder policy}

Every final module should satisfy:
\begin{enumerate}[label=(\roman*),leftmargin=2.3em]
\item no \texttt{sorry};
\item no theorem-specific \texttt{axiom};
\item public endpoints are literal \texttt{Prf PRAAx};
\item \texttt{\#print axioms} shows only the metatheoretic axioms already accepted by the repository audit;
\item every semantic test is labeled separately from object derivation status.
\end{enumerate}
The blueprint files in this archive intentionally violate the first condition by visible milestones; they are not presented as finished modules.

\section{Theorem ledger and exact status}
\label{sec:status}

\small
\begin{longtable}{@{}p{0.28\textwidth}p{0.45\textwidth}p{0.19\textwidth}@{}}
\toprule
Result & Evidence in this archive/repository & Status \\
\midrule
\endfirsthead
\toprule
Result & Evidence in this archive/repository & Status \\
\midrule
\endhead
Ordinary \(\Prf\;\PRAAx\) calculus, substitution, weakening, cut & Existing \path{ConciseIPRA/Deduction.lean} & \textcolor{statusgreen}{\textbf{checked}} \\
Open quantifier-free induction and instantiation & Existing \path{IPRADerivations/Core/Basic.lean} & \textcolor{statusgreen}{\textbf{checked}} \\
Bounded least search and minimality & Existing \path{IPRADerivations/Core/Search.lean} & \textcolor{statusgreen}{\textbf{checked}} \\
Pairing, equality bits, simultaneous recursion & Existing \path{Core/Pairing.lean} & \textcolor{statusgreen}{\textbf{checked}} \\
Finite sums/products and primitive-recursive definitions & Existing \path{Core/Definitions.lean} & \textcolor{statusgreen}{\textbf{checked}} \\
Fixed bounded Boolean reflection & Full human proof, Lean blueprint & \human{} / \outline \\
Selected master certificate & Human reduction to checked search API & \human{} \\
Deterministic trace reflection & Full human proof, direct \texttt{ind} route & \human{} \\
Finite monotone demand closure & Full rank and leastness proof; tested reference implementation & \human \\
Backward viability and least successful path & Full finite proof; tested reference implementation & \human \\
Exact/quotient/metric observer erasure & Full human proof & \human \\
Finite observational descent & Full tree-induction proof; tested metric example & \human \\
Residual closure & Direct wrapper around existing \texttt{Prf.trans} & \textcolor{statusgreen}{\textbf{checked rule}} \\
Generic FINER compiler & Human metatheorem and source-oriented blueprint & \human{} / \outline \\
Python algorithms & Fourteen exact tests pass & Semantic test only \\
Cubic bisection benchmark & Exact Python certificate and developed object-proof plan & Semantic + \outline \\
General Faltings/Siegel/Mordell--Weil producers & Not supplied & \openstatus \\
Uniform exact Puiseux/root/eigenvector selector through collisions & Deliberately not claimed; interface reformulated & Obstructed \\
General rate extraction from metastability & Deliberately not claimed & Obstructed \\
Open-sequent Foundation bridge & Repository unification-plan work, not implemented here & Future integration \\
\bottomrule
\end{longtable}
\normalsize

\subsection{Files in the archive}

\begin{description}[style=nextline,leftmargin=2.5em]
\item[\texttt{paper/}] Complete LaTeX source and compiled PDF.
\item[\texttt{lean/FINER/}] A tiny kernel-wrapper file plus source-oriented theorem blueprints.  The blueprints visibly contain \texttt{sorry} and are not claimed compiled.
\item[\texttt{lean/FORMALIZATION\_PLAN.md}] Dependency-ordered implementation checklist.
\item[\texttt{reference/}] Exact standard-library Python algorithms and unit tests.
\item[\texttt{examples/}] JSON fixtures for demand closure, relational viability, metric descent, and cubic bisection.
\item[\texttt{metadata/}] Theorem ledger, difficulty matrix, schema, proof-status record, and source audit.
\item[\texttt{validation/}] Test output, PDF build log, preflight information, rendered pages, and static scans.
\end{description}

\subsection{Reproduction commands}

From the archive root:
\begin{lstlisting}[language=bash]
python3 -m unittest -v reference/test_finer_reference.py
./build.sh
\end{lstlisting}
The build script runs tests, compiles the PDF, checks for missing references, renders the PDF, and refreshes checksums.  Lean validation requires placing \texttt{lean/} in the pinned ProveIt checkout and following the formalization plan; the current archive does not assert that the blueprint modules elaborate.

\section{Conclusion and recommended next steps}
\label{sec:conclusion}

The existing ProveIt development already contains the logical kernel and much of the arithmetic needed for a broad finitary mathematics program.  The main integration problem is to stop treating every theorem packet as a new foundational experiment.  FINER proposes one stable compilation boundary:
\begin{center}
finite query \(\longrightarrow\) least demand closure \(\longrightarrow\) one checked master certificate \(\longrightarrow\) ordinary \(\Prf\;\PRAAx\) endpoint with explicit residuals.
\end{center}

The framework resolves the shared engineering and proof-theoretic bottlenecks: bounded logic is reflected to bits; loops use one quantifier-free induction; finite choice is replaced by viability; witness dependence is erased only through exact, quotient, or metric observers; cyclic demands have a least finite solution; and missing quantitative mathematics remains visible in the context.  It also gives principled negative answers: no universal primitive-recursive evaluator, no unguarded exact branch selector through collisions, no general global rate from finite metastability, and no height cutoff from cardinality alone.

The highest-value implementation sequence is now clear.

\begin{enumerate}[label=\textbf{Step \arabic*.},leftmargin=4.2em]
\item Check a compact Boolean-reflection library against the current \texttt{IPRADerivations} API.
\item Formalize deterministic trace reflection and the fixed-dimensional demand-closure theorem.
\item Formalize fixed-horizon viability and exact observer erasure.
\item Package selected-certificate soundness and residual closure in a small \texttt{CompiledQuery} record.
\item Replay the four benchmarks in \cref{sec:lean} with literal \texttt{Prf PRAAx} endpoints.
\item Migrate one existing proof packet and publish a kernel-checked theorem ledger entry.
\end{enumerate}

This route is deliberately incremental.  Every completed phase creates useful checked lemmas even if the full graph DSL is never built.  Conversely, a large source-language compiler without these small object-theory kernels would merely add another unverified layer.  The framework therefore prioritizes a narrow trusted path and exact theorem status over breadth of front-end notation.


\appendix
\section{Current ProveIt API map}
\label{app:api}

This appendix records the concrete repository declarations on which the formalization plan is based.  It is an integration map, not a substitute for checking the pinned sources.

\begin{longtable}{@{}p{0.29\textwidth}p{0.27\textwidth}p{0.36\textwidth}@{}}
\toprule
Need & Declaration/module & Intended use \\
\midrule
\endfirsthead
\toprule
Need & Declaration/module & Intended use \\
\midrule
\endhead
Ordinary derivability & \texttt{ConciseIPRA.Prf} & Public object-theory endpoint \\
Substitution and weakening & \texttt{Prf.subst}, \texttt{Prf.weaken} & Specialize reusable local proofs \\
Simultaneous cut & \texttt{Prf.trans} & Compose graph nodes and close residuals \\
Closed theorem reuse & \texttt{Prf.ofNil} & Use shared arithmetic under local context \\
Open induction & \texttt{LeanProofs.IPRA.ind} & Trace, fold, rank, and table invariants \\
Closed derivation instantiation & \texttt{LeanProofs.IPRA.inst} & Arbitrary terms and contexts \\
Primitive-recursive term compiler & \texttt{PTm.toPR}, \texttt{toPR\_eq} & Turn fixed terms into symbols \\
Primitive recursion & \texttt{recDef}, \texttt{rec1}, \texttt{rec2} & States, folds, viability tables \\
Finite sums/products & \texttt{sumF}, \texttt{prodF} & Boolean folds, ranks, error accumulation \\
Equality bit & \texttt{eqb}, \texttt{eqb\_of\_eq}, \texttt{eqb\_of\_ne} & Atomic reflection and branching \\
Pairing/projections & \texttt{pairT}, \texttt{fstT}, \texttt{sndT} & Packed states and certificates \\
Simultaneous recursion & \texttt{Core.Pairing} simultaneous section & Multi-component state machines \\
Order & \texttt{Core.Order} & Bounds, monotonicity, finite minima \\
Bounded search & \texttt{Search.mu} & Least accepted certificate \\
Search specification & \texttt{search\_spec} & Bound, success, minimality in one theorem \\
Known witness search & \texttt{search\_found} & Derive selected certificate from existence \\
Minimality & \texttt{search\_minimal} & No earlier accepted code \\
Division/remainder & \texttt{Core.Division} & Mixed-radix digits and arrays \\
Finite multi-sums & \texttt{Core.MultiSums} & Variable-size tables and error sums \\
Ring normalization & \texttt{Core.Tactic.ipra\_ring} & Polynomial identities such as cubic monotonicity \\
Foundation export & \texttt{IPRABridge} closed interfaces & Backend corollary after ordinary proof \\
\bottomrule
\end{longtable}

\subsection{Representative exact theorem shapes}

The checked open induction wrapper has the conceptual shape
\begin{lstlisting}[language={}]
theorem ind (phi : Fm PTm (k+1)) (hqf : phi.QF)
    (h0 : [] |- [PRAAx] phi[0])
    (hS : [phi] |- [PRAAx] phi[S(var 0)]) :
    [] |- [PRAAx] phi
\end{lstlisting}
with precise repository substitution notation.

The checked search theorem has the conceptual consequence
\begin{lstlisting}[language={}]
search_found
  (hf : Gamma |- [PRAAx] f(i,x) = 0)
  (hi : Gamma |- [PRAAx] i < N) :
  Gamma |- [PRAAx]
    mu(f,N,x) < N /\
    f(mu(f,N,x),x) = 0 /\
    mu(f,N,x) <= i
\end{lstlisting}
This is exactly the selected-certificate lemma after defining \(f=1\dot{-}\CertCheck\).

The checked simultaneous-cut theorem has the conceptual shape
\begin{lstlisting}[language={}]
theorem Prf.trans
    (h : Gamma |- [Ax] phi)
    (s : forall psi in Gamma, Delta |- [Ax] psi) :
    Delta |- [Ax] phi
\end{lstlisting}
which is the public residual-discharge operation.

\subsection{Import-boundary recommendation}

The generic FINER library should have a narrow dependency boundary: the ordinary calculus and only the small arithmetic modules it actually uses.  Application modules may import the bridge or Foundation, but the finite compiler kernels should not import Busy Beaver, ZFC-in-PA, Diophantine certificates, or packet-specific data.  This follows the repository's unification plan and keeps compilation costs and theorem provenance understandable.

\section{Executable reference algorithms}
\label{app:reference}

The Python module \path{reference/finer_reference.py} uses only the standard library and \path{fractions.Fraction}.  It is intentionally independent of Lean and contains no symbolic theorem prover.  Its role is to make the finite algorithms and certificate formats executable before object formalization.

\subsection{Implemented algorithms}

\begin{description}[style=nextline,leftmargin=2.7em]
\item[Mixed-radix coding.] \texttt{mixed\_radix\_encode}, \texttt{mixed\_radix\_decode}, and \texttt{certificate\_space\_size}.

\item[Least accepted search.] \texttt{least\_accepted} returns the least accepted code or the bound, and \texttt{verify\_search\_result} checks success or exhaustive failure.

\item[Demand closure.] \texttt{monotone\_demand\_closure} computes the join iteration, and \texttt{verify\_demand\_closure} checks every step and the final prefix condition.  \texttt{verify\_least\_prefixed} brute-forces leastness for small examples.

\item[Backward viability.] \texttt{backward\_viability} computes viable layers and the least viable path; \texttt{verify\_viability} checks the complete dynamic-programming table and path.

\item[Observers.] Exact invariance, finite diameter, and spanning-tree error accumulation are implemented for finite values.

\item[Height enumeration.] \texttt{bounded\_complete\_list} and its verifier illustrate the height-bound to list direction.

\item[Cubic bisection.] Exact rational bisection for \(y^3+y=x\) returns a replayable certificate.  The verifier recomputes the complete trace and final bracket.
\end{description}

\subsection{Test coverage}

The fourteen unit tests cover:
\begin{enumerate}[label=(\arabic*),leftmargin=2.3em]
\item mixed-radix round trip;
\item overflow rejection;
\item successful least search;
\item exhaustive failed search;
\item cyclic three-coordinate demand closure;
\item leastness of a nontrivial two-coordinate closure;
\item a viable relational path;
\item a nonviable relational instance;
\item exact observer invariance;
\item metric observer diameter;
\item spanning-tree metric gluing;
\item bounded complete-list verification;
\item cubic bisection;
\item rejection of a tampered cubic certificate.
\end{enumerate}
The delivered run reports all tests passing.  Test output is stored in \texttt{validation/test\_output.txt}.

\subsection{Relationship to the object proof}

The reference functions mirror the primitive-recursive algorithms intended for Lean.  They do not establish that the corresponding \texttt{PTm} computes the same function or that the object theory proves the advertised contract.  Each formalization milestone therefore has three separate checks:
\begin{enumerate}[label=(\roman*),leftmargin=2.3em]
\item semantic tests of the reference algorithm;
\item definitional/evaluation comparison between the Lean \texttt{PR} program and the reference on fixtures;
\item a kernel-checked \texttt{Prf PRAAx} correctness theorem.
\end{enumerate}
Only the third changes theorem status.

\section{Source audit and relationship to prior packets}
\label{app:audit}

\subsection{Pinned repository state}

The report was prepared against ProveIt main commit
\begin{center}
\texttt{64bc0e997df6b8dbdf43f5ac9b3d9f74b4047ee5}
\end{center}
on 16 November 2026.  The audit focused on:
\begin{itemize}[leftmargin=2em]
\item \texttt{Logic/PrimitiveRecursiveArithmetic/README.md};
\item \texttt{UNIFICATION-PLAN.md};
\item \texttt{ConciseIPRA/README.md} and ordinary calculus sources;
\item \texttt{Derivations/README.md} and core arithmetic sources;
\item \texttt{Bridge/README.md};
\item the foundations packets, especially CFP, finite-use, FQNF, CBIP, KIFC-R, and their difficulty matrices.
\end{itemize}

\subsection{Inherited ideas}

FINER deliberately inherits and consolidates the following ideas rather than claiming them as new:
\begin{itemize}[leftmargin=2em]
\item finite-query specialization;
\item certified finitary presentations and explicit moduli;
\item bounded Boolean reflection;
\item master certificates and least search;
\item reverse demand propagation;
\item proof schedules composed by simultaneous cut;
\item relational witness clouds and backward viability;
\item stable or quotient observer erasure;
\item producer/verifier separation;
\item residual endpoints of the form \(B_1,\ldots,B_r,\Gamma\vdash\Phi\).
\end{itemize}
These themes occur in several packets.  The repository's current README correctly marks their new Lean files as uncompiled outlines or candidates unless a separate checked derivation exists.

\subsection{Added synthesis}

This report contributes the following sharper integration choices.

\begin{enumerate}[label=(\arabic*),leftmargin=2.3em]
\item The ordinary \texttt{Prf PRAAx} interface is fixed as the sole canonical endpoint; native or Foundation exports are downstream.
\item Cyclic demands are handled by a fully stated least-prefixed-point theorem on finite products of chains, with an explicit \(\sum C_i\) rank bound and overflow discipline.
\item Finite observational descent is isolated as a reusable theorem that avoids constructing global quotients or witnesses when only an invariant observation is needed.
\item The selected-certificate theorem is tied directly to the already checked \texttt{Core.Search} API, avoiding duplicate search developments.
\item The formalization plan begins with small fixed-dimension and fixed-horizon kernels, postponing a graph DSL until object derivations are checked.
\item The theorem ledger treats external semantic evidence as a third status, neither derived nor residual, so it cannot silently enter the compiled proof.
\end{enumerate}

\subsection{Nonduplication with the unification plan}

The repository's unification plan addresses transport between ordinary ConciseIPRA and Foundation, sharing arithmetic, and migrating application consumers.  FINER sits above the ordinary arithmetic library and below applications.  It does not replace the bridge plan.  In particular:
\begin{itemize}[leftmargin=2em]
\item formula and open-sequent Foundation transport remain bridge work;
\item symbol-comparison lemmas remain shared-arithmetic work;
\item FINER produces ordinary object proofs that those layers can later transport;
\item application-specific tables, polynomial coefficients, and consistency predicates remain outside the generic compiler.
\end{itemize}

\subsection{Audit caveat}

The repository was changing rapidly on 14--16 November 2026.  This archive pins the commit above and names exact source paths.  Later repository changes may rename declarations or complete some blueprints.  Such changes should update the source audit and theorem ledger rather than retroactively altering the status claimed here.


\begin{thebibliography}{20}
\bibitem{ProveItRoot}
V.~Reshetnikov, \emph{ProveIt}, directory \texttt{Logic/PrimitiveRecursiveArithmetic}, repository snapshot \texttt{64bc0e997d}, 16 November 2026.

\bibitem{ConciseIPRA}
V.~Reshetnikov, \emph{Concise presentations of intuitionistic primitive recursive arithmetic}, ProveIt/Logic/PrimitiveRecursiveArithmetic/ConciseIPRA.

\bibitem{Derivations}
V.~Reshetnikov, \emph{iPRA derivations of theorems from the proof packets}, ProveIt/Logic/PrimitiveRecursiveArithmetic/Derivations.

\bibitem{Bridge}
V.~Reshetnikov, \emph{The Foundation--ConciseIPRA bridge}, ProveIt/Logic/PrimitiveRecursiveArithmetic/Bridge.

\bibitem{Unification}
V.~Reshetnikov, \emph{ProveIt iPRA unification plan}, \texttt{UNIFICATION-PLAN.md}.

\bibitem{FQNF}
ProveIt research packet, \emph{Finite-Query Normal Forms for Mathematics in iPRA}, 16 November 2026.

\bibitem{KIFCR}
ProveIt research packet, \emph{KIFC-R: residualized finitary compilation for iPRA}, 16 November 2026.

\bibitem{CBIP}
ProveIt research packet, \emph{Certified Bounded Interaction Programs for iPRA}, 16 November 2026.

\bibitem{FiniteUse}
ProveIt research packet, \emph{Demand-Driven Relational Certificates for Mathematics in iPRA}, 15 November 2026.

\bibitem{CFP}
ProveIt research packet, \emph{Certified Finitary Presentations for Mathematics in iPRA}, 15 November 2026.
\end{thebibliography}

\end{document}
